\chapter{Elliptic boundary problems}\label{ch:boundary}
\providecommand{\dCmI}{\dot{\mathcal{C}}^{-\infty}}
\providecommand{\Hdot}[1]{\dot{H}^{#1}}
\providecommand{\Xt}{\tilde{X}}
\providecommand{\PiC}{\Pi_{\mathcal{C}}}
\providecommand{\Cau}{\mathfrak{C}}
\providecommand{\Rat}[1]{\mathcal{R}^{#1}}
\providecommand{\cJ}{\mathcal{J}}
\providecommand{\Fin}{\mathcal{F}}

An elliptic differential operator on a compact manifold \emph{with boundary} has, in
general, an infinite-dimensional null space. The basic problem of the theory of elliptic
boundary problems is to parametrize this null space by boundary data, and then to
decide which boundary conditions single out a Fredholm problem. In this chapter we
carry this out, following the approach of Calder\'on, for first-order operators of Dirac
type. The analysis rests on three ingredients: a careful description of function spaces
and distributions on a manifold with boundary, culminating in the \emph{jump formula}
(an invariant form of Green's formula); the observation that a manifold with boundary
can always be regarded as a closed subset of a manifold without boundary, so that the
calculus of \cref{ch:manifolds} is available; and an explicit analysis of the boundary
behaviour of $Q(w\,\delta(x))$, where $Q$ is a parametrix for a Dirac-type operator and
$w\,\delta(x)$ is a layer supported on a hypersurface. For Dirac-type operators the
symbol of $Q$ is rational in the fibre variables, and contour deformation shows that
these layer potentials are smooth up to the hypersurface from each side, with boundary
values given by pseudodifferential operators.

We treat first the model situation of a hypersurface $H$ dividing a compact manifold
$M$ without boundary, on which $D$ is assumed invertible; here the Calder\'on
projector $\PiC$, a pseudodifferential projection on $H$ whose range is the space of
boundary values of solutions on one side, is obtained by an explicit formula. We then
treat a compact manifold with boundary, where the same construction works modulo
smoothing errors, and where some abstract functional analysis of smoothing operators
and generalized inverses is needed to correct them. Finally we show that a
pseudodifferential projection $\Pi$ on the boundary defines a Fredholm boundary
problem $\{u;\ \Pi(u|_{\pa X})=0\}$ precisely when $\Pi$ restricted to the range of
$\PiC$ is Fredholm, and give a symbolic criterion for this.

Throughout, $\Psido{m}(Y;E,F)$ for a compact manifold $Y$ without boundary is as in
\cref{ch:manifolds}; ``classical'' means that the local full symbols are
polyhomogeneous (in $\Symph{m}$), and then the principal symbol $\sigma_m$ is a
homogeneous function on $T^*Y\setminus0$. A manifold is always assumed paracompact and
Hausdorff, and vector bundles are complex.

\section{Manifolds with boundary}\label{boundary:sec:manifolds}

The local model for a manifold with boundary is the half-space
\begin{equation*}
\R^{n,1}=[0,\infty)\times\R^{n-1}.
\end{equation*}
An \emph{open set} in $\R^{n,1}$ means a relatively open set, $O=O'\cap\R^{n,1}$ with
$O'\subset\R^n$ open. A map $F:O_1\to O_2$ between open subsets of $\R^{n,1}$ is
\emph{smooth} if it is the restriction of a smooth map: there are open sets
$O_i'\subset\R^n$ with $O_i=O_i'\cap\R^{n,1}$ and a $\CI$ map $F':O_1'\to O_2'$ with
$F=F'|_{O_1}$. This is much stronger than smoothness of $F$ on $O_1\cap\{x_1>0\}$. A
\emph{diffeomorphism} is a smooth bijection with smooth inverse.

\begin{definition}
A \emph{coordinate system} on a topological space $X$ is a homeomorphism
$\Phi:U\to O$ from an open set $U\subset X$ onto an open set $O\subset\R^{n,1}$. Two
coordinate systems $(U_i,\Phi_i,O_i)$ are \emph{compatible} if $U_1\cap U_2=\emptyset$
or $\Phi_2\circ\Phi_1^{-1}:\Phi_1(U_1\cap U_2)\to\Phi_2(U_1\cap U_2)$ is a
diffeomorphism. A \emph{$\CI$ manifold with boundary} is a paracompact Hausdorff space
with a maximal covering by mutually compatible coordinate systems.
\end{definition}

Equivalently, $X$ is a manifold with boundary if there is a manifold $\Xt$ without
boundary and $f\in\CI(\Xt)$, real-valued, with $df\ne0$ on $\{f=0\}$, such that
$X=\{p\in\Xt;\ f(p)\ge0\}$ with the coordinate systems obtained by restriction. One
direction is clear; the doubling construction below gives the other.

The \emph{boundary} $\pa X$ is the union of the sets $\Phi^{-1}(O\cap(\{0\}\times
\R^{n-1}))$ over coordinate systems; it is well defined because a diffeomorphism of
open subsets of $\R^{n,1}$ maps boundary points to boundary points (its extension is a
local diffeomorphism of $\R^n$, hence an open map, so it cannot map a boundary point to
an interior point of $\R^{n,1}$ and still map a neighbourhood into $\R^{n,1}$, and
symmetrically for the inverse). $\pa X$ is a manifold without boundary, compact if $X$
is, and $X^\circ=X\setminus\pa X$ is a manifold without boundary.

A \emph{boundary defining function} is a real $\rho\in\CI(X)$ (see the next section)
with $\rho\ge0$, $\pa X=\{\rho=0\}$ and $d\rho\ne0$ on $\pa X$. Such functions exist:
locally the first coordinate is one, and these may be patched together with a
partition of unity, since positive combinations of local defining functions are again
defining functions.

\begin{lemma}[Collar neighbourhoods]\label{boundary:lem:collar}
Let $X$ be a manifold with compact boundary and $\rho$ a boundary defining function.
Then there exist $\epsilon>0$, an open neighbourhood $C$ of $\pa X$ and a
diffeomorphism $\varphi:C\to[0,\epsilon)\times\pa X$ whose first component is $\rho$.
\end{lemma}
\begin{proof}
Choose a Riemannian metric on $X$ (by patching local ones) and let $V$ be a smooth
vector field equal to $\nabla\rho/|\nabla\rho|^2$ near $\pa X$; it is inward pointing
and $V\rho=1$ near $\pa X$. For $p\in\pa X$ the integral curve $t\mapsto\exp(tV)p$
exists for $0\le t<\epsilon$, for some $\epsilon>0$ uniform by compactness of $\pa X$,
and $\rho(\exp(tV)p)=t$ for small $t$. The map $\psi:[0,\epsilon)\times\pa X\ni(t,p)
\mapsto\exp(tV)p$ is smooth, its differential at $t=0$ is an isomorphism (it maps
$\pa_t$ to $V$, transversal to $\pa X$), and it is injective since $\rho(\psi(t,p))=t$
and the curves starting at distinct boundary points are distinct. Shrinking $\epsilon$,
$\psi$ is a diffeomorphism onto an open neighbourhood $C$; take $\varphi=\psi^{-1}$.
\end{proof}

Using a collar we can \emph{double} $X$. Let
\begin{equation*}
\Xt=(X\sqcup X)/\pa X
\end{equation*}
be two copies of $X$ with corresponding boundary points identified, and define
$\CI(\Xt)$ to consist of the pairs $(u_1,u_2)\in\CI(X)\oplus\CI(X)$ for which there
is $f\in\CI((-\epsilon,\epsilon)\times\pa X)$ with $(\varphi^{-1})^*(u_1|_C)=
f(\rho,\cdot)$ and $(\varphi^{-1})^*(u_2|_C)=f(-\rho,\cdot)$. This is a $\CI$ structure
on $\Xt$, a manifold without boundary which is compact if $X$ is, the first copy
embeds $X$ as $\{x\ge0\}$ where $x\in\CI(\Xt)$ is equal to $\rho$ on the first copy
and $-\rho$ on the second near $\pa X$, and $\CI(X)=\CI(\Xt)|_X$. Vector bundles over
$X$ double in the same way (using the collar to identify $E$ near $\pa X$ with the
pull-back of $E|_{\pa X}$). Thus all the usual objects --- tensor bundles, densities,
differential operators --- may be obtained by restriction from $\Xt$.

\section{Smooth functions}\label{boundary:sec:smooth}

For $O\subset\R^{n,1}$ open there are two natural definitions of smooth functions:
\begin{equation}\label{boundary:eq:CI-ext}
\CI(O)=\{u:O\to\C;\ u=\tilde u|_O\ \text{for some}\ \tilde u\in\CI(O'),\ O'\subset
\R^n\ \text{open},\ O=O'\cap\R^{n,1}\}
\end{equation}
and the intrinsic one: $u\in\CI(O^\circ)$, $O^\circ=O\cap((0,\infty)\times\R^{n-1})$,
with
\begin{equation}\label{boundary:eq:CI-int}
\sup_{K\cap O^\circ}|D^\alpha u|<\infty\quad\text{for all}\ K\Subset O\ \text{and
all}\ \alpha.
\end{equation}
Certainly \eqref{boundary:eq:CI-ext} implies \eqref{boundary:eq:CI-int}. Conversely,
\eqref{boundary:eq:CI-int} implies that each $D^\alpha u$ is uniformly continuous on
$K\cap O^\circ$, hence extends continuously to $O$, and these extensions are the
derivatives of the extension of $u$ in the closed half-space.

\begin{proposition}\label{boundary:prop:extension}
If $u$ satisfies \eqref{boundary:eq:CI-int}, then $u$ extends to a smooth function on
an open set $O'\subset\R^n$ with $O'\cap\R^{n,1}=O$; hence the two definitions agree.
Moreover there is a linear extension map $E:\CI(O)\to\CI(O')$, continuous in the sense
that for every compact $K'\subset O'$ and $k$ there is a compact $K\subset O$ with
$\sup_{K'}|D^\alpha Eu|\le C\sum_{|\beta|\le k}\sup_{K}|D^\beta u|$ for $|\alpha|\le k$.
\end{proposition}
\begin{proof}
We prove the existence of extensions, which is all that is used below. By a
partition of unity we may assume $u$ has compact support in $O$. Let
$u_k(y)=\pa_x^ku(0,y)$, which are in $\CIc(\R^{n-1})$. By Borel's lemma with
parameters (the usual proof, $\tilde u(x,y)=\sum_k\chi(x/\epsilon_k)x^ku_k(y)/k!$ with
$\epsilon_k\downarrow0$ chosen so that the $k$th term is bounded by $2^{-k}$ in
$\mathcal{C}^{k-1}$, works verbatim; see \cite[Section~1.2]{Hormander-I}) there is
$\tilde u\in\CIc(\R^n)$ with $\pa_x^k\tilde u(0,y)=u_k(y)$ for all $k$. Set $Eu=u$ for
$x\ge0$ and $Eu=\tilde u$ for $x<0$. All derivatives of $Eu$ from the two sides are
continuous and agree on $x=0$, so $Eu\in\mathcal{C}^k$ for every $k$ (a function
which is $\mathcal{C}^1$ on each closed half-space with matching derivatives is
$\mathcal{C}^1$, by the mean value theorem, and one iterates). A linear continuous
extension with the stated bounds is Seeley's extension operator $Eu(x,y)=
\sum_k a_k\psi(-b_kx)u(-b_kx,y)$ for $x<0$, with $b_k=2^k$, $\psi$ a cut-off and
$a_k$ chosen with $\sum_ka_k(-b_k)^j=1$ for all $j$ and $\sum_k|a_k|b_k^j<\infty$; see
the Notes.
\end{proof}

For a manifold $X$ with boundary, $\CI(X)$ is the space of functions whose
coordinate representations are smooth in this sense; it suffices to check this for
one covering by coordinate systems. As for manifolds without boundary, $\CI(X)$ admits
partitions of unity, and its topology is given by the seminorms $\sup_K|D^\alpha
(\Phi^{-1})^*(u|_U)|$, $K\Subset O$, of which a countable subset suffices; for $X$
compact $\CI(X)$ is a Fr\'echet space. The restriction map $\CI(X)\ni u\mapsto
u|_{\pa X}\in\CI(\pa X)$ is surjective: using a collar, extend $v\in\CI(\pa X)$ to be
independent of $\rho$ and cut off.

We think of $\CI(X)\subset\CI(X^\circ)$; a differential operator $P\in\Diff{m}(X)$ is
a differential operator on $X^\circ$ with coefficients smooth up to the boundary, and
$P:\CI(X)\to\CI(X)$. The space of functions vanishing to infinite order at the
boundary is
\begin{equation*}
\dCI(X)=\{u\in\CI(X);\ (Pu)|_{\pa X}=0\ \ \forall\,P\in\Diff{*}(X)\}.
\end{equation*}
Extension by zero identifies $\dCI(X)$ with the closed subspace of $\CI(\Xt)$ of
functions supported in $X$. Given a collar with normal variable $x=\rho$, Taylor's
theorem and Borel's lemma give an isomorphism
\begin{equation*}
\CI(X)/\dCI(X)\cong\CI(\pa X)[[x]],\quad u\mapsto\sum_k\frac{x^k}{k!}\,
(\pa_x^ku)|_{\pa X},
\end{equation*}
onto formal power series with coefficients in $\CI(\pa X)$. The same applies to
sections of vector bundles.

\section{Distributions}\label{boundary:sec:distributions}

Let $X$ be a compact manifold with boundary and $\Xt$ a compact manifold without
boundary containing $X$ as $\{x\ge0\}$ (for example the double). There are three
spaces of distributions associated with $X$. The coarsest is $\CmI(X^\circ)$, the dual
of $\CIc(X^\circ;\Omega)$; its elements may behave arbitrarily badly at $\pa X$ and
play only a conceptual role. The direct analogue of the boundaryless case is the space
of \emph{supported distributions}
\begin{equation*}
\dCmI(X)=\big(\CI(X;\Omega)\big)'.
\end{equation*}
Since the restriction map $\CI(\Xt;\Omega)\to\CI(X;\Omega)$ is continuous and
surjective, duality gives an injective map
\begin{equation}\label{boundary:eq:supp-incl}
\dCmI(X)\ni u\longmapsto\tilde u\in\CmI(\Xt),\quad \tilde u(\varphi)=u(\varphi|_X).
\end{equation}

\begin{proposition}\label{boundary:prop:supported}
The range of \eqref{boundary:eq:supp-incl} is $\{\tilde u\in\CmI(\Xt);\
\supp\tilde u\subset X\}$.
\end{proposition}

The proof is given after \cref{boundary:prop:extendible}. We use
\eqref{boundary:eq:supp-incl} to identify $\dCmI(X)$ with the distributions on $\Xt$
supported in $X$; this is consistent with the identification of $\dCI(X)$ with
smooth functions on $\Xt$ supported in $X$. The closed subspace $\dCI(X;\Omega)\subset
\CI(\Xt;\Omega)$ has dual the space of \emph{extendible distributions}
\begin{equation*}
\CmI(X)=\big(\dCI(X;\Omega)\big)',
\end{equation*}
and transposing the inclusion gives the restriction map
\begin{equation}\label{boundary:eq:RX}
R_X:\CmI(\Xt)\to\CmI(X),\quad R_Xu(\varphi)=u(\varphi),\ \varphi\in\dCI(X;\Omega).
\end{equation}
This is not an identification: its null space is
$\{v\in\CmI(\Xt);\ \supp v\cap X^\circ=\emptyset\}=\dCmI(\Xt\setminus X^\circ)$, the
distributions supported ``on the other side of the boundary''.

\begin{proposition}\label{boundary:prop:extendible}
$\CIc(X^\circ;\Omega)$ is dense in $\dCI(X;\Omega)$. Hence restriction gives an
injective map $\CmI(X)\hookrightarrow\CmI(X^\circ)$, while $R_X:\dCmI(X)\to\CmI(X)$
(and a fortiori $R_X$ on $\CmI(\Xt)$) is surjective.
\end{proposition}
\begin{proof}
Let $V$ be a real vector field on $\Xt$ which is inward pointing across $\pa X$, and
$F_s=\exp(sV)$, a one-parameter group of diffeomorphisms of $\Xt$ with $F_s(X)\subset
X^\circ$ for $s>0$. If $\varphi\in\dCI(X)$ then $F_s^*\varphi\to\varphi$ in $\CI(\Xt)$
as $s\to0$, and for $s<0$, $F_s^*\varphi=\varphi\circ F_s$ is supported in
$F_{-s}(X)\Subset X^\circ$, so lies in $\CIc(X^\circ)$. The same argument applies to
densities. Injectivity of $\CmI(X)\to\CmI(X^\circ)$ follows by duality. Surjectivity
of $R_X$ on $\dCmI(X)$ follows from the Hahn--Banach theorem: a continuous linear
functional on the closed subspace $\dCI(X;\Omega)\subset\CI(X;\Omega)$ extends to
$\CI(X;\Omega)$.
\end{proof}

\begin{proof}[Proof of \cref{boundary:prop:supported}]
For $\tilde u\in\CmI(\Xt)$ the condition $\supp\tilde u\subset X$ means $\tilde
u(\varphi)=0$ for all $\varphi\in\CIc(\Xt\setminus X;\Omega)$. This certainly holds
if $\tilde u$ comes from $\dCmI(X)$, since then $\varphi|_X=0$. Conversely, if it
holds, then by continuity and the density of $\CIc(\Xt\setminus X;\Omega)$ in
$\dCI(\Xt\setminus X^\circ;\Omega)$ (\cref{boundary:prop:extendible} applied to the
manifold with boundary $\Xt\setminus X^\circ$), $\tilde u$ vanishes on the null space
of the restriction $\CI(\Xt;\Omega)\to\CI(X;\Omega)$, so it factors through this
(open, by the open mapping theorem) surjection.
\end{proof}

On $\dCmI(X)$ and $\CmI(X)$ we may use the weak topologies, given by the seminorms
$u\mapsto|u(\varphi)|$ for test functions $\varphi$.

\begin{lemma}\label{boundary:lem:weak-density}
$\CIc(X^\circ)$ is weakly dense in both $\dCmI(X)$ and $\CmI(X)$.
\end{lemma}
\begin{proof}
For $u\in\dCmI(X)$, with $F_s$ as above and $s>0$, the push-forward $u_s=(F_s)_*u$,
$u_s(\varphi)=u(F_s^*\varphi)$, has compact support in $X^\circ$, and $u_s\to u$
weakly since $F_s^*\varphi\to\varphi$ in $\CI$. A distribution with compact support
in $X^\circ$ is the weak limit of elements of $\CIc(X^\circ)$ (mollify in local
coordinates after a partition of unity). Weak density being transitive, the first
claim follows. For $u\in\CmI(X)$ choose by \cref{boundary:prop:extendible}
$\tilde u\in\CmI(\Xt)$ with $R_X\tilde u=u$; $\tilde u$ is a weak limit of
$\psi_k\in\CI(\Xt)$ (on a compact manifold without boundary), so $u$ is the weak limit
of $\psi_k|_X\in\CI(X)$. Finally $\psi\in\CI(X)$ is the weak limit in $\CmI(X)$ of
$\chi_k\psi$ with $\chi_k\in\CIc(X^\circ)$, $0\le\chi_k\le1$, $\chi_k\to1$ pointwise
on $X^\circ$, by dominated convergence.
\end{proof}

\section{Boundary terms and the jump formula}\label{boundary:sec:jump}

To relate supported and extendible distributions we use the space of
\emph{boundary terms}
\begin{equation*}
\dCmI_{\pa X}(X)=\{u\in\dCmI(X);\ \supp u\subset\pa X\},
\end{equation*}
the support being computed in $\Xt$, and the map ``cutting off at the boundary''
\begin{equation*}
\CI(X)\ni u\longmapsto u_c\in\dCmI(X),\quad u_c(\varphi)=\int_Xu\varphi\ \ \forall\,
\varphi\in\CI(X;\Omega),
\end{equation*}
which in terms of $\Xt$ is extension by zero. The following is immediate from the
definitions and \cref{boundary:prop:extendible}.

\begin{proposition}\label{boundary:prop:exact}
For a compact manifold with boundary $X$ there is a commutative diagram
\begin{equation*}
\begin{tikzcd}
& & & \dCI(X)\arrow[d,hook] & \\
& & & \CI(X)\arrow[d,hook]\arrow[dl,"(\cdot)_c"'] & \\
0\arrow[r] & \dCmI_{\pa X}(X)\arrow[r] & \dCmI(X)\arrow[r,"R_X"] & \CmI(X)\arrow[r] & 0
\end{tikzcd}
\end{equation*}
with exact bottom row.
\end{proposition}
\begin{proof}
$R_X(u_c)=u$ for $u\in\CI(X)$, which is the commutativity. Surjectivity of $R_X$ is
in \cref{boundary:prop:extendible}. If $v\in\dCmI_{\pa X}(X)$ then $v(\varphi)=0$ for
$\varphi\in\CIc(X^\circ;\Omega)$, hence by density for $\varphi\in\dCI(X;\Omega)$, so
$R_Xv=0$. Conversely if $u\in\dCmI(X)$ and $R_Xu=0$ then $u(\varphi)=0$ for
$\varphi\in\CIc(X^\circ;\Omega)$, so $\supp u\cap X^\circ=\emptyset$ and $\supp u
\subset\pa X$.
\end{proof}

Both maps from $\CI(X)$ into distributions are injective; we regard the map into
$\CmI(X)$ as an identification, but \emph{not} the map $u\mapsto u_c$.

\subsection*{Actions of differential operators}
Let $E,F$ be vector bundles over $X$ and $P\in\Diff{m}(X;E,F)$, with transpose
$P^t\in\Diff{m}(X;F^*\otimes\Omega,E^*\otimes\Omega)$ defined by $\int_X\pair{P^t\varphi}{u}=
\int_X\pair{\varphi}{Pu}$ for $u\in\CIc(X^\circ;E)$, $\varphi\in\CIc(X^\circ;F^*\otimes
\Omega)$. Integration by parts shows that this identity continues to hold if
$u\in\CI(X;E)$ and $\varphi\in\dCI(X;F^*\otimes\Omega)$, or vice versa, but
\emph{not} for $u,\varphi$ both merely smooth up to the boundary. Accordingly $P$ acts
in two ways on distributions, by duality:
\begin{align*}
P&:\CmI(X;E)\to\CmI(X;F), & (Pu)(\varphi)&=u(P^t\varphi),\ \varphi\in\dCI(X;F^*\otimes
\Omega),\\
\dot P&:\dCmI(X;E)\to\dCmI(X;F), & (\dot Pu)(\varphi)&=u(P^t\varphi),\
\varphi\in\CI(X;F^*\otimes\Omega).
\end{align*}
The first extends the action on $\CI(X;E)$. If $\tilde P$ is any extension of $P$ to a
differential operator on $\Xt$, then $\dot P$ is the restriction of $\tilde P$ to
distributions supported in $X$, since $\tilde P^t\tilde\varphi$ restricted to $X$
depends only on $\tilde\varphi|_X$. The failure of $\dot P(u_c)=(Pu)_c$ for $u\in
\CI(X;E)$ is measured by a boundary term, $\dot P(u_c)-(Pu)_c\in\dCmI_{\pa X}(X;F)$,
which we now compute: this is Green's formula in distributional form.

Fix a collar $C\cong[0,\epsilon)_x\times\pa X$ as in \cref{boundary:lem:collar}, and
identify $E$, $F$ over $C$ with the pull-backs of $E|_{\pa X}$, $F|_{\pa X}$ (possible
since $C$ retracts onto $\pa X$). On the double, $x$ extends to
$(-\epsilon,\epsilon)\times\pa X$ and we have the distributions
\begin{equation*}
v\,D_x^j\delta(x)\in\dCmI_{\pa X}(X;F),\quad v\in\CI(\pa X;F),\ j\ge0,
\end{equation*}
defined in local coordinates $(x,y)$ as $v(y)\otimes D_x^j\delta(x)$, i.e.\ by
$(v\,D^j_x\delta(x))(\varphi\,|dx\,dy|)=\int v(y)\cdot((-D_x)^j\varphi)(0,y)\,dy$;
this is independent of the choice of $y$. With $H(x)$ the Heaviside function, $u_c=
H(x)u$ near $\pa X$, and $\delta(x)=\pa_xH(x)=iD_xH(x)$.

\begin{proposition}[Jump formula]\label{boundary:prop:jump}
For each $P\in\Diff{m}(X;E,F)$ there are differential operators $P_{jk}\in
\Diff{m-1-j-k}(\pa X;E,F)$, $j,k\ge0$, $j+k\le m-1$, such that for all $u\in\CI(X;E)$
\begin{equation}\label{boundary:eq:jump-general}
\dot P(u_c)-(Pu)_c=\sum_{j+k\le m-1}\big(P_{jk}\,(D_x^ku)|_{\pa X}\big)\,D_x^j\delta(x),
\end{equation}
and $P_{0,m-1}=-i\,\sigma_m(P)(dx)|_{\pa X}$.
\end{proposition}
\begin{proof}
Both sides vanish if $u\in\dCI(X;E)$ (then $u_c$ is smooth on $\Xt$), so only the
Taylor series of $u$ at $\pa X$ matters and we may work in $C$. There
\begin{equation*}
P=\sum_{l=0}^mP_l(x)D_x^l,
\end{equation*}
where $P_l(x)\in\Diff{m-l}(\pa X;E,F)$ depends smoothly on $x$, and
$P_m(0)=\sigma_m(P)(dx)$ (with the convention that $D_x$ has symbol $\xi$). From
$D_x(Hf)=HD_xf-i\delta(x)f(0,\cdot)$ one finds by induction
\begin{equation*}
D_x^l(Hu)=HD_x^lu-i\sum_{k=0}^{l-1}\big((D_x^ku)|_{\pa X}\big)\,D_x^{l-1-k}\delta(x).
\end{equation*}
The tangential operators $P_l(x)$ commute with multiplication by $H$. Expanding
$P_l(x)=\sum_{r<m}x^rP_{l,r}+x^mP_l'(x)$ by Taylor's formula and using $x^rD^p_x\delta
(x)=c_{r,p}D_x^{p-r}\delta(x)$, where $c_{r,p}=i^rp!/(p-r)!$ for $r\le p$ and
$c_{r,p}=0$ for $r>p$ (in particular $c_{0,p}=1$), we obtain
\eqref{boundary:eq:jump-general} with
\begin{equation*}
P_{jk}=-i\sum_{l,r:\ l-1-k-r=j}c_{r,l-1-k}P_{l,r}.
\end{equation*}
Here $P_{l,r}$ has order at most $m-l\le m-1-j-k$. The only contribution to
$P_{0,m-1}$ comes from $l=m$, $r=0$, giving $-iP_m(0)=-i\sigma_m(P)(dx)$.
\end{proof}

For first-order operators there is no choice involved:

\begin{corollary}\label{boundary:cor:jump1}
If $D\in\Diff{1}(X;E,F)$ and $u\in\CI(X;E)$ then
\begin{equation}\label{boundary:eq:jump}
\dot D(u_c)=(Du)_c+\frac1i\,\sigma_1(D)(dx)\big(u|_{\pa X}\big)\,\delta(x).
\end{equation}
\end{corollary}

The right side does not depend on the choice of $x$: if $x'=ax$ with $a>0$ then on
$\pa X$, $\sigma_1(D)(dx')=a\sigma_1(D)(dx)$ while $\delta(x')=a^{-1}\delta(x)$. In
the sequel we write $b$ for the restriction map, $bu=u|_{\pa X}$.

The same discussion applies to an embedded hypersurface $H$ with a defining function
$x$ in a manifold without boundary: if $u$ is smooth on each of the two sides
$\{\pm x\ge0\}$ near $H$ (i.e.\ $u|_{\{\pm x>0\}}$ extends smoothly to $\{\pm x\ge0\}$),
with boundary values $b_\pm u$, then applying \eqref{boundary:eq:jump} on both sides,
\begin{equation}\label{boundary:eq:jump-two}
Du=(Du)_{+,c}+(Du)_{-,c}+\frac1i\sigma_1(D)(dx)\big(b_+u-b_-u\big)\delta(x),
\end{equation}
where $(Du)_{\pm,c}$ is $Du$ computed in $\{\pm x>0\}$, extended by zero.

\subsection*{Formal solutions at the boundary}
For an elliptic first-order operator, the Cauchy problem can always be solved
formally at the boundary.

\begin{lemma}\label{boundary:lem:formal}
Let $D\in\Diff{1}(X;V_1,V_2)$ with $\sigma_1(D)(dx)$ invertible at $\pa X$ (this holds
if $D$ is elliptic).
\begin{enumerate}[label=(\roman*)]
\item For $f\in\CI(X;V_2)$ and $v\in\CI(\pa X;V_1)$ there is $u\in\CI(X;V_1)$ with
$bu=v$ and $Du-f\in\dCI(X;V_2)$.
\item If $u\in\CI(X;V_1)$, $bu=0$ and $Du\in\dCI(X;V_2)$, then $u\in\dCI(X;V_1)$.
\end{enumerate}
\end{lemma}
\begin{proof}
In the collar, shrinking it if necessary, $D=A(x)\big(D_x+E(x)\big)$ with $A(x)$ an
invertible bundle map ($A(0)=\sigma_1(D)(dx)$) and $E(x)\in\Diff{1}(\pa X;V_1)$
depending smoothly on $x$. Write $u_k=(D_x^ku)|_{\pa X}$. The condition that
$D_x^k\big((D_x+E)u-A^{-1}f\big)$ vanish at $x=0$ for all $k$ is
\begin{equation*}
u_{k+1}=\big(D_x^k(A^{-1}f)\big)|_{x=0}-\sum_{j\le k}\binom kj\big(D_x^{k-j}E\big)(0)\,u_j,
\quad k\ge0.
\end{equation*}
For (i) set $u_0=v$, define $u_k$ recursively and use Borel's lemma to find $u$ with
these Taylor coefficients. For (ii), the recursion with $f$ replaced by $Du$ (all of
whose Taylor coefficients vanish) and $u_0=0$ gives $u_k=0$ for all $k$.
\end{proof}

\section{Sobolev spaces}\label{boundary:sec:sobolev}

As for smooth functions, the (extendible) Sobolev spaces on a compact manifold with
boundary $X\subset\Xt$ are defined by restriction,
\begin{equation*}
H^m(X)=R_X\big(H^m(\Xt)\big)\subset\CmI(X),\quad \norm{u}_{H^m(X)}=\inf\{\norm{\tilde
u}_{H^m(\Xt)};\ R_X\tilde u=u\},\quad m\in\R,
\end{equation*}
the quotient norm; independence of the choice of $\Xt$ follows from the locality of
Sobolev spaces and their invariance under diffeomorphisms. The \emph{supported}
Sobolev spaces are
\begin{equation*}
\Hdot{m}(X)=\{u\in H^m(\Xt);\ \supp u\subset X\}\subset\dCmI(X),
\end{equation*}
closed subspaces of $H^m(\Xt)$. Both are $\CI(X)$-modules, and for a vector bundle $E$
one sets $H^m(X;E)=H^m(X)\otimes_{\CI(X)}\CI(X;E)$, and similarly for $\Hdot{m}$.

\begin{lemma}\label{boundary:lem:Hdot-dense}
$\CIc(X^\circ)$ is dense in $\Hdot{m}(X)$ for every $m\in\R$.
\end{lemma}
\begin{proof}
With $F_s$ as in the proof of \cref{boundary:prop:extendible}, the operators $u\mapsto
(F_s)_*u$ are uniformly bounded on $H^m(\Xt)$ for $|s|\le1$ (their norms depend on
finitely many derivatives of $F_s$) and $(F_s)_*u\to u$ for $u\in\CI(\Xt)$, hence for
all $u\in H^m(\Xt)$ by density. For $u\in\Hdot{m}(X)$ and $s>0$, $(F_s)_*u\in
H^m(\Xt)$ has compact support in $X^\circ$, and is approximated in $H^m$ by
mollification in local coordinates.
\end{proof}

\begin{proposition}\label{boundary:prop:duality}
For $m\in\R$ the pairing $(u,v)\mapsto\int_Xuv$, $u\in\CI(X)$, $v\in\dCI(X;\Omega)$,
extends by continuity to a non-degenerate pairing identifying $H^m(X)$ with the dual
of $\Hdot{-m}(X;\Omega)$, and $\Hdot{m}(X)$ with the dual of $H^{-m}(X;\Omega)$.
\end{proposition}
\begin{proof}
The dual of the closed subspace $\Hdot{-m}(X;\Omega)\subset H^{-m}(\Xt;\Omega)$ is the
quotient of $H^m(\Xt)=(H^{-m}(\Xt;\Omega))'$ by the annihilator $\{v\in H^m(\Xt);\
v(w)=0\ \forall\,w\in\Hdot{-m}(X;\Omega)\}$. By \cref{boundary:lem:Hdot-dense} it
suffices to test with $w\in\CIc(X^\circ;\Omega)$, so the annihilator is $\{v;\
v|_{X^\circ}=0\}$, the null space of $R_X$ on $H^m(\Xt)$. Hence the dual is
$H^m(\Xt)/\Nul(R_X)\cong H^m(X)$ with the quotient norm. The second statement follows
since Hilbert spaces are reflexive.
\end{proof}

A differential operator $P\in\Diff{k}(X;E,F)$ defines bounded maps
\begin{equation*}
P:H^m(X;E)\to H^{m-k}(X;F),\qquad \dot P:\Hdot{m}(X;E)\to\Hdot{m-k}(X;F),
\end{equation*}
by restriction from an extension $\tilde P$, consistently with the actions on
distributions. Moreover
\begin{equation*}
\CmI(X;E)=\bigcup_mH^m(X;E),\quad \CI(X;E)=\bigcap_mH^m(X;E),\quad
\dCmI(X;E)=\bigcup_m\Hdot{m}(X;E),\quad \dCI(X;E)=\bigcap_m\Hdot{m}(X;E).
\end{equation*}
For the second: if $u\in H^m(X)$ for all $m$, then for each $k$, $u$ is the
restriction of a $\mathcal{C}^k$ function on $\Xt$ (Sobolev embedding), so all
derivatives of $u$ are bounded on $X^\circ$ and $u\in\CI(X)$ by
\cref{boundary:prop:extension}. The others follow from the corresponding statements
on $\Xt$.

\subsection*{Layers on a hypersurface}
We shall need the Sobolev regularity of distributions supported on a hypersurface.
In $\R^n=\R_x\times\R^{n-1}_y$, for $v\in\CmIc(\R^{n-1})$ the distribution $v\otimes
\delta(x)$ has Fourier transform $\hat v(\eta)$, so, since
$\int_\R(1+|\eta|^2+\xi^2)^t\,d\xi=c_t\ang{\eta}^{2t+1}$ with $c_t=\int(1+s^2)^tds
<\infty$ precisely for $t<-\frac12$,
\begin{equation}\label{boundary:eq:layer-norm}
\norm{v\otimes\delta(x)}^2_{H^t(\R^n)}=\frac{c_t}{2\pi}\norm{v}^2_{H^{t+\frac12}(\R^{n-1})},
\quad t<-\tfrac12.
\end{equation}

\begin{lemma}\label{boundary:lem:layers}
Let $H\subset M$ be a compact embedded hypersurface, with defining function $x$ near
$H$, in a manifold $M$ without boundary, and $E$ a vector bundle over $M$.
\begin{enumerate}[label=(\roman*)]
\item For $t<-\frac12$, $v\mapsto v\,\delta(x)$ is bounded and injective
$H^{t+\frac12}(H;E)\to H^t_{\mathrm{c}}(M;E)$, with closed range; in particular
for $t=-1$ it is an isomorphism of $H^{-\frac12}(H;E)$ onto a closed subspace of
$H^{-1}_{\mathrm c}(M;E)$.
\item Every $T\in H^{-1}_{\mathrm{c}}(M;E)$ with $\supp T\subset H$ is of the form
$T=v\,\delta(x)$ with $v\in H^{-\frac12}(H;E)$.
\end{enumerate}
\end{lemma}
\begin{proof}
(i) follows from \eqref{boundary:eq:layer-norm} by a partition of unity, which also
shows that the norms of $v$ and $v\,\delta(x)$ are equivalent. For (ii), localize: we
may assume $T\in H^{-1}(\R^n)$ has compact support in $\{x=0\}$. If $\varphi\in
\CIc(\R^n)$ vanishes on $\{x=0\}$, let $\chi\in\CIc(\R)$ equal $1$ near $0$ and
$\varphi_\epsilon=(1-\chi(x/\epsilon))\varphi$. Then $T(\varphi_\epsilon)=0$, and
$\varphi_\epsilon\to\varphi$ in $H^1$: indeed $\chi(x/\epsilon)\varphi$ is supported
in $|x|\le C\epsilon$ where $|\varphi|\le C|x|$, so its $L^2$ norm is $O(\epsilon^{3/2})$
and that of its derivative is $O(\epsilon^{1/2})$ (the term $\epsilon^{-1}\chi'(x/\epsilon)
\varphi$ being bounded). Hence $T(\varphi)=0$, i.e.\ $T(\varphi)$ depends only on
$\varphi(0,\cdot)$, and $T=v\otimes\delta(x)$ where $v(\psi)=T(\chi(x)\psi(y))$ is a
compactly supported distribution. Then $\hat T(\xi,\eta)=\hat v(\eta)$ and
$\norm{T}_{H^{-1}}^2=c\int|\hat v(\eta)|^2\ang{\eta}^{-1}d\eta<\infty$, so $v\in
H^{-\frac12}$.
\end{proof}

\section{Generalized Dirac operators}\label{boundary:sec:dirac}

For $D\in\Diff{1}(M;V_1,V_2)$, in local coordinates and trivializations
$D=\sum_jA_j(z)D_{z_j}+A'(z)$, and the principal symbol is
$\sigma_1(D)(z,\zeta)=\sum_jA_j(z)\zeta_j$, linear in $\zeta\in T^*_zM$.

\begin{definition}\label{boundary:def:dirac}
Let $M$ be a manifold with or without boundary, with a Riemannian metric $g$, and let
$V_1,V_2$ be Hermitian vector bundles over $M$. An operator $D\in\Diff{1}(M;V_1,V_2)$
is a \emph{generalized Dirac operator} (or \emph{of Dirac type}) if
\begin{equation}\label{boundary:eq:dirac}
\sigma_1(D)(z,\zeta)^*\,\sigma_1(D)(z,\zeta)=|\zeta|_g^2\,\Id_{V_1}\quad\text{for all}
\ (z,\zeta)\in T^*M,
\end{equation}
and $\operatorname{rank}V_1=\operatorname{rank}V_2$. If $M_+\subset M$ is a subset,
$D$ is of Dirac type on $M_+$ if \eqref{boundary:eq:dirac} holds for $z\in M_+$.
\end{definition}

We abbreviate $\sigma(\zeta)=\sigma_1(D)(z,\zeta)$ and write $\tau(\zeta)=
\sigma(\zeta)^*$ for real $\zeta$, extended complex-linearly to complex covectors, and
$g(\zeta)=|\zeta|^2_g$, extended as a complex quadratic form. Then $\sigma(\zeta)$ is
invertible for $\zeta\ne0$ with
\begin{equation*}
\sigma(\zeta)^{-1}=\frac{\tau(\zeta)}{g(\zeta)},\qquad \sigma(\zeta)\tau(\zeta)=
g(\zeta)\Id_{V_2},
\end{equation*}
since a left inverse of a square matrix is a right inverse. Thus $D$ is elliptic, its
formal adjoint $D^*$ (with respect to the Riemannian density and the Hermitian
structures) is again of Dirac type, with $\sigma_1(D^*)=\sigma_1(D)^*$, and $D^*D$,
$DD^*$ are second order with principal symbol $|\zeta|^2_g\Id$. Polarizing
\eqref{boundary:eq:dirac} and its counterpart on $V_2$ gives the \emph{Clifford
relations}
\begin{equation}\label{boundary:eq:clifford}
\sigma(\zeta)^*\sigma(\zeta')+\sigma(\zeta')^*\sigma(\zeta)=2g(\zeta,\zeta')\Id_{V_1},\quad
\sigma(\zeta)\sigma(\zeta')^*+\sigma(\zeta')\sigma(\zeta)^*=2g(\zeta,\zeta')\Id_{V_2}
\end{equation}
for real covectors $\zeta,\zeta'$.

\begin{example}\label{boundary:ex:dirac}
(1) The Cauchy--Riemann operator $D=D_x+iD_y$ on $\R^2$ (or on a Riemann surface with
a compatible metric, acting on functions with values in $(0,1)$-forms) has
$\sigma(\xi,\eta)=\xi+i\eta$, $\sigma^*\sigma=\xi^2+\eta^2$.
(2) On a Riemannian manifold, $d+d^*$ acting from even to odd forms has symbol
$i(\zeta\wedge-\iota_\zeta)$, and $(\zeta\wedge-\iota_\zeta)^2=-|\zeta|^2$; so it is
of Dirac type, as is $d+d^*$ on all forms.
(3) The Dirac operator on the spinor bundle of a spin manifold, possibly twisted by a
Hermitian bundle with connection, is of Dirac type; this is the origin of the name.
\end{example}

\section{Dividing hypersurfaces}\label{boundary:sec:dividing}

To isolate the essential idea of Calder\'on we first consider the following
situation.

\medskip\noindent\textbf{Standing hypotheses (dividing case).} $M$ is a compact
Riemannian manifold without boundary, $D\in\Diff{1}(M;V_1,V_2)$ is a generalized
Dirac operator, and $x\in\CI(M)$ is real with $dx\ne0$ on $H=\{x=0\}$; thus $H$ is a
compact hypersurface with trivial normal bundle, and
\begin{equation*}
M=M_+\cup M_-,\quad M_\pm=\{\pm x\ge0\},\quad M_+\cap M_-=H,
\end{equation*}
where $M_\pm$ are compact manifolds with boundary $H$. In addition we assume
\begin{equation}\label{boundary:eq:invertible}
D:\CI(M;V_1)\to\CI(M;V_2)\ \text{is an isomorphism.}
\end{equation}

By \cref{ch:manifolds}, $D$ is Fredholm, so \eqref{boundary:eq:invertible} forces the
index to vanish; it is made to simplify the discussion. By the results there,
$D^{-1}\in\Psido{-1}(M;V_2,V_1)$, and $D^{-1}$ extends to an isomorphism
$\CmI(M;V_2)\to\CmI(M;V_1)$ inverting $D$ on distributions.

Our object is to parametrize the spaces of solutions
\begin{equation*}
\mathcal{N}_\pm=\{u\in\CI(M_\pm;V_1);\ Du=0\ \text{in}\ M_\pm^\circ\}
\end{equation*}
by their boundary values. Let $V=V_1|_H$ and $b_\pm:\CI(M_\pm;V_1)\to\CI(H;V)$ the
restriction maps.

\begin{lemma}\label{boundary:lem:unique}
Under the standing hypotheses $b_\pm:\mathcal{N}_\pm\to\CI(H;V)$ are injective.
\end{lemma}
\begin{proof}
If $u\in\mathcal{N}_+$ and $b_+u=0$, then by \cref{boundary:lem:formal}(ii) $u$
vanishes to infinite order at $H$. Hence $\tilde u$, equal to $u$ in $M_+$ and to $0$
in $M_-$, is smooth on $M$, and $D\tilde u=0$. By \eqref{boundary:eq:invertible},
$\tilde u=0$.
\end{proof}

The main result in this setting is the following.

\begin{theorem}[Calder\'on projector, dividing case]\label{boundary:thm:calderon-dividing}
Under the standing hypotheses there is a unique $\PiC\in\Psido{0}(H;V)$ with
$\PiC^2=\PiC$ such that
\begin{equation*}
b_+:\mathcal{N}_+\to\PiC\,\CI(H;V)\quad\text{and}\quad
b_-:\mathcal{N}_-\to(\Id-\PiC)\,\CI(H;V)
\end{equation*}
are isomorphisms. In particular $\CI(H;V)=b_+(\mathcal{N}_+)\oplus b_-(\mathcal{N}_-)$.
$\PiC$ is classical and, if $\nu=dx/|dx|$ is the unit conormal and $T^*H$ is identified
with the covectors orthogonal to $\nu$,
\begin{equation}\label{boundary:eq:calderon-symbol}
\sigma_0(\PiC)(y,\eta)=\frac12\Big(\Id+\frac{i\,\sigma(\nu)^{-1}\sigma(\eta)}{|\eta|}
\Big),\quad(y,\eta)\in T^*H\setminus0,
\end{equation}
the orthogonal projection onto the $+1$ eigenspace of the self-adjoint involution
$i\sigma(\nu)^{-1}\sigma(\eta)/|\eta|$.
\end{theorem}

The construction uses the jump formula ``in reverse''. Consider the space of
piecewise solutions
\begin{equation*}
\mathcal{N}_H=\{u\in L^1(M;V_1);\ u|_{M_\pm^\circ}\in\mathcal{N}_\pm\}\cong
\mathcal{N}_+\oplus\mathcal{N}_-.
\end{equation*}
By \eqref{boundary:eq:jump-two}, for $u\in\mathcal{N}_H$,
\begin{equation}\label{boundary:eq:jump-solutions}
Du=\frac1i\sigma(dx)\big(b_+u-b_-u\big)\,\delta(x)\in\CI(H;V_2)\,\delta(x).
\end{equation}

\begin{proposition}\label{boundary:prop:layer-iso}
Under the standing hypotheses, $D:\mathcal{N}_H\to\CI(H;V_2)\,\delta(x)$ is an
isomorphism with inverse $w\,\delta(x)\mapsto D^{-1}(w\,\delta(x))$.
\end{proposition}
\begin{proof}
Injectivity: if $u\in\mathcal{N}_H$ and $Du=0$ on $M$ then $u=D^{-1}Du=0$. For
surjectivity we must show that for $w\in\CI(H;V_2)$,
\begin{equation}\label{boundary:eq:layer-smooth}
u=D^{-1}(w\,\delta(x))\ \text{satisfies}\ u|_{M_\pm^\circ}\in\CI(M_\pm;V_1);
\end{equation}
then $Du=0$ in $M_\pm^\circ$ and $u\in\mathcal{N}_H$ (note $u\in H^{1/2-\epsilon}
\subset L^1$ by \cref{boundary:lem:layers}). This is the case $Q=D^{-1}$ of
\cref{boundary:thm:layer} below, applied on both sides of $H$.
\end{proof}

The regularity statement \eqref{boundary:eq:layer-smooth} is the analytic heart of the
matter. It is \emph{false} for a general pseudodifferential operator of order $-1$ in
place of $D^{-1}$ (for instance for $\Op(\ang{\zeta}^{-1})$ in $\R^n$, $n\ge2$, the
result is not smooth up to $H$); it depends on the special form of the symbol of
$D^{-1}$, which we now examine.

\section{Rational symbols and layer potentials}\label{boundary:sec:rational}

\subsection*{Symbols over a closed set}
Let $W\subset\R^n$ be open and $K\subset W$ closed in $\R^n$ (we shall take $K$ a
compact subset of a closed half-space). We say $a\in\CI(W\times\R^n)$ is a symbol of
order $m$ \emph{over $K$}, $a\in S^m(K)$, if for all $\alpha,\beta$,
$\sup_{z\in K}|D_z^\alpha D_\zeta^\beta a(z,\zeta)|\le C_{\alpha\beta}\ang{\zeta}^{m-|\beta|}$.
These classes are closed under products, $D_z$ and $D_\zeta$ in the usual way. For a
symbol $q\in S^m(\R^n;\R^n)$, $\Op(q)f(z)$ depends only on $q(z,\cdot)$.

Let $D=\sum_jA_j(z)D_{z_j}+A'(z)$ be a first-order operator on $W$ (with matrix
coefficients), with full symbol $d(z,\zeta)=\sigma(z,\zeta)+A'(z)$. Since $d$ is
linear in $\zeta$, for $q\in S^m(\R^n;\R^n)$ with $z$-support in $W$,
\begin{equation}\label{boundary:eq:comp-diff}
D\,\Op(q)=\Op(d\circ q),\qquad d\circ q=dq+\sum_jA_jD_{z_j}q,
\end{equation}
exactly, as follows by differentiating under the integral sign.

\subsection*{Rational symbols}
Suppose now that $W$ carries a Riemannian metric with dual metric $g(z,\zeta)=
\sum g^{jk}(z)\zeta_j\zeta_k$ and Hermitian structures, and that $\tau(z,\zeta)$ is
linear in $\zeta$, with $\sigma\tau=g$ and $\tau\sigma=g$ for $z\in K$. (If $D$ is of
Dirac type on $K$, take $\tau=\sigma^*$ computed with respect to the Hermitian
structures; outside $K$ the identities need not hold.)

\begin{definition}
For $r\in\Z$, $\Rat{r}$ is the space of finite sums of matrix-valued functions
$p(z,\zeta)/(1+g(z,\zeta))^j$, where $j\ge0$ and $p$ is a polynomial in $\zeta$ of
degree at most $2j+r$ with coefficients in $\CI(W)$.
\end{definition}

Clearly $\Rat{r}\cdot\Rat{s}\subset\Rat{r+s}$, $\Rat{r}\subset\Rat{r+1}$, and
$\Rat{r}$ is preserved by $D_{z_j}$ and multiplication by $\CI(W)$; moreover
$\pa_{\zeta_j}\Rat{r}\subset\Rat{r-1}$, since $\pa_{\zeta_j}(p(1+g)^{-j})=
(\pa_{\zeta_j}p)(1+g)^{-j}-jp(\pa_{\zeta_j}g)(1+g)^{-j-1}$. Since $1+g\ge c\ang{\zeta}^2$
on compact sets, it follows that $\Rat{r}\subset S^r(L)$ for every compact $L\subset W$.

\begin{lemma}[Rational symbols]\label{boundary:lem:rational}
In this setting there are $\rho_N\in\Rat{-1}$, $N\ge1$, with $\rho_1=\tau/(1+g)$ and
$\rho_{N+1}-\rho_N\in\Rat{-N-1}$, such that the following holds. If $L\subset K$ is
compact and $q\in S^{-1}(\R^n;\R^n)$ satisfies $d\circ q-1\in S^{-\infty}(L')$ for a
compact neighbourhood $L'$ of $L$ in $K$, then
\begin{equation*}
q-\rho_N\in S^{-N-1}(L)\quad\text{for all}\ N.
\end{equation*}
Similarly, with $\phi\in\CIc(\R^n)$ equal to $1$ near $0$,
$q-\sum_{l=1}^N(1-\phi)p_lg^{1-2l}\in S^{-N-1}(L)$, where $p_1=\tau$ and
\begin{equation*}
p_{l+1}=-\tau\Big(gA'p_l+\sum_jA_j\big(gD_{z_j}p_l-(2l-1)(D_{z_j}g)p_l\big)\Big)
\end{equation*}
are polynomials in $\zeta$, homogeneous of degree $3l-2$. Thus the homogeneous terms of
the symbol of a parametrix of a generalized Dirac operator are of the form
$g^{1-2l}p_l$.
\end{lemma}
\begin{proof}
Over $K$, $d\circ\rho_1=g/(1+g)+A'\tau/(1+g)+\sum_jA_jD_{z_j}(\tau/(1+g))=1-e_1$ with
$e_1\in\Rat{-1}$. Inductively, if $d\circ\rho_k=1-e_k$ over $K$ with $e_k\in
\Rat{-k}$, set $\rho_{k+1}=\rho_k+\tau e_k/(1+g)$. Then $\tau e_k/(1+g)\in\Rat{-k-1}$
and over $K$
\begin{equation*}
d\circ\frac{\tau e_k}{1+g}=e_k-\frac{e_k}{1+g}+A'\frac{\tau e_k}{1+g}+\sum_jA_j
D_{z_j}\frac{\tau e_k}{1+g},
\end{equation*}
so $d\circ\rho_{k+1}=1-e_{k+1}$ with $e_{k+1}\in\Rat{-k-1}$. (The identity
$\sigma\tau=g$ is used only pointwise; the derivatives fall on the smooth functions
$\tau,g$.)

Now let $e=q-\rho_N\in S^{-1}(L')$. Then $d\circ e=(d\circ q-1)+e_N\in S^{-N}(L')$.
Suppose $e\in S^{-k}(L'')$ for a compact neighbourhood $L''\subset L'$ of $L$, $k\le
N$. Then $\sigma e=d\circ e-A'e-\sum_jA_jD_{z_j}e\in S^{-k}(L'')$ and
$e=(1-\phi)\tau g^{-1}\sigma e+\phi e\in S^{-k-1}(L'')$. Starting from $k=1$ we
reach $e\in S^{-N-1}(L)$ (the neighbourhoods may be fixed in advance, since only
derivatives at points of $L''$ are involved). The homogeneous version is proved in
the same way: writing $q_{-l}=p_lg^{1-2l}$, the terms of order $-l$ in
$d\circ\sum_lq_{-l}=1$ give $\sigma q_{-1}=1$ and $q_{-l-1}=-g^{-1}\tau\big(A'q_{-l}+
\sum_jA_jD_{z_j}q_{-l}\big)$, and substituting
$D_{z_j}q_{-l}=g^{-2l}\big(gD_{z_j}p_l-(2l-1)(D_{z_j}g)p_l\big)$ gives the recursion for
$p_l$; the degree of $p_{l+1}$ is $1+2+(3l-2)=3(l+1)-2$.
\end{proof}

\subsection*{The basic contour integral}
Let $W=(-1,1)_x\times O$, $O\subset\R^{n-1}$ open, so $\zeta=(\xi,\eta)$, and write
\begin{equation*}
g(z,\xi,\eta)=g^{11}(z)\xi^2+2\beta(z,\eta)\xi+g'(z,\eta),\quad
\beta=\sum_{k\ge2}g^{1k}\eta_k.
\end{equation*}
By positive-definiteness $g^{11}>0$ and $g^{11}g'(\eta)-\beta(\eta)^2\ge c|\eta|^2$
locally uniformly. Hence for each $(z,\eta)$ the equation $1+g(z,\xi,\eta)=0$ has the
two roots
\begin{equation*}
\xi_\pm(z,\eta)=\frac{-\beta(z,\eta)\pm i\,s(z,\eta)}{g^{11}(z)},\quad
s=\big(g^{11}(1+g')-\beta^2\big)^{1/2}\ge c\ang{\eta},
\end{equation*}
which lie in the upper and lower half planes, at distance $\ge c\ang{\eta}$ from the
real axis and from each other, with $|\xi_\pm|\le C\ang{\eta}$. For $\rho=p/(1+g)^j\in
\Rat{r}$ with $r\le-1$ and $x\ge0$ set
\begin{equation}\label{boundary:eq:Gplus}
G^+_\rho(z',x,\eta)=\frac1{2\pi}\oint_{\gamma}e^{ix\xi}\rho(z',\xi,\eta)\,d\xi=
i\operatorname*{Res}_{\xi=\xi_+(z',\eta)}\big(e^{ix\xi}\rho(z',\xi,\eta)\big),
\end{equation}
where $\gamma$ is a positively oriented circle in $\{\operatorname{Im}\xi>0\}$
enclosing $\xi_+(z',\eta)$ but not $\xi_-(z',\eta)$; similarly
$G^-_\rho(z',x,\eta)=-i\operatorname{Res}_{\xi=\xi_-}(e^{ix\xi}\rho)$ for $x\le0$.

\begin{lemma}\label{boundary:lem:contour}
Let $\rho\in\Rat{r}$, $r\le-1$, and let $L\subset W$ be compact.
\begin{enumerate}[label=(\roman*)]
\item For fixed $(z',\eta)$ the inverse Fourier transform of $\xi\mapsto\rho(z',\xi,
\eta)$ (an $L^2$ function) is $H(x)G^+_\rho(z',x,\eta)-H(-x)G^-_\rho(z',x,\eta)$.
\item For $z'\in L$, $x\ge0$ and all $\alpha,k,\beta$,
$|D^\alpha_{z'}D^k_xD^\beta_\eta G^+_\rho(z',x,\eta)|\le C\ang{\eta}^{r+1+k-|\beta|}$.
In particular $G^+_\rho(z',0,\eta)\in S^{r+1}(L)$ (as a symbol in $\eta$).
\item For $z'\in L$, $G^+_\rho(z',0,\eta)$ is classical: for $|\eta|$ large it has an
asymptotic expansion in functions homogeneous in $\eta$ of degrees $r+1-l$, $l\ge0$.
\item If $\rho=\tau/(1+g)$, the term of degree $0$ in (iii) is
$i\,\tau(z',\xi^0_+,\eta)/\big(g^{11}(\xi^0_+-\xi^0_-)\big)$, where $\xi^0_\pm=
(-\beta\pm i(g^{11}g'-\beta^2)^{1/2})/g^{11}$ are the roots of $g(z',\xi,\eta)=0$.
\end{enumerate}
\end{lemma}
\begin{proof}
(i) Since $r<0$, $\rho$ is a proper rational function of $\xi$ and has a partial
fraction decomposition $\sum_{k=1}^j\big(a_k(\xi-\xi_+)^{-k}+b_k(\xi-\xi_-)^{-k}\big)$.
From $\int_0^\infty e^{-ix(\xi-\xi_+)}dx=-i(\xi-\xi_+)^{-1}$ (convergent as
$\operatorname{Im}\xi_+>0$) we get $\cF^{-1}[(\xi-\xi_+)^{-1}]=iH(x)e^{ix\xi_+}$, and
differentiating in $\xi_+$, $\cF^{-1}[(\xi-\xi_+)^{-k}]=iH(x)(ix)^{k-1}e^{ix\xi_+}/
(k-1)!=iH(x)\operatorname{Res}_{\xi_+}\big(e^{ix\xi}(\xi-\xi_+)^{-k}\big)$; similarly
$\cF^{-1}[(\xi-\xi_-)^{-k}]=-iH(-x)\operatorname{Res}_{\xi_-}(e^{ix\xi}(\xi-
\xi_-)^{-k})$. Summing gives (i).

(ii) Near a given $(z_0,\eta_0)$ we may fix $\gamma$ to be the circle with centre
$\xi_+(z_0,\eta_0)$ and radius $c'\ang{\eta_0}$, $c'$ small; it encloses
$\xi_+(z',\eta)$ and lies in $\operatorname{Im}\xi\ge c''\ang{\eta}$ for $(z',\eta)$
in a neighbourhood of size comparable to $\ang{\eta_0}$ in $\eta$ (by the estimates on
$\xi_\pm$ and $|\pa_{z,\eta}\xi_\pm|\le C$, $|\pa_z\xi_\pm|\le C\ang\eta$). On $\gamma$,
$|\xi|\le C\ang{\eta}$ and $|1+g|=g^{11}|\xi-\xi_+||\xi-\xi_-|\ge c\ang{\eta}^2$, so
$|\rho|\le C\ang{\eta}^r$, and likewise $|D_{z'}^\alpha D_\eta^\beta\rho|\le
C\ang\eta^{r-|\beta|}$ because $D_z^\alpha D_\eta^\beta\rho\in\Rat{r-|\beta|}$.
Moreover $|e^{ix\xi}|\le1$ and $|D_x^ke^{ix\xi}|\le C\ang{\eta}^k$ for $x\ge0$. As
$\gamma$ has length $O(\ang{\eta})$, differentiation under the integral gives (ii).

(iii) For $|\eta|\ge1$ write $\eta=|\eta|\theta$, $\xi=|\eta|t$ and $\epsilon=
|\eta|^{-2}$; then $1+g=|\eta|^2(\epsilon+g(z',t,\theta))$ and, decomposing $p$ into
homogeneous parts $p_{(k)}$,
\begin{equation*}
G^+_\rho(z',0,\eta)=\sum_k|\eta|^{k+1-2j}\,\frac1{2\pi}\oint_{\gamma'}\frac{p_{(k)}(z',t,
\theta)}{(\epsilon+g(z',t,\theta))^j}\,dt,
\end{equation*}
where $\gamma'$ is a fixed circle around the root $t^0_+(z',\theta)$ of
$g(z',t,\theta)=0$ in the upper half plane, which for $0\le\epsilon\le\epsilon_0$ still
encloses exactly one root of $\epsilon+g=0$. The integrals are smooth in
$(z',\theta,\epsilon)$ for $0\le\epsilon\le\epsilon_0$; Taylor expansion in $\epsilon=
|\eta|^{-2}$ gives an expansion in homogeneous functions of degrees $k+1-2j-2l\le
r+1-2l$.

(iv) This is the term $k=1$, $l=0$, i.e.\ $\epsilon=0$: the residue of $\tau/g$ at the
simple pole $\xi^0_+$.
\end{proof}

\subsection*{Boundary behaviour of layer potentials}

\begin{theorem}\label{boundary:thm:layer}
Let $M$ be a compact manifold without boundary, $x\in\CI(M)$ real with $dx\ne0$ on
$H=\{x=0\}$, $M_+=\{x\ge0\}$, and let $D\in\Diff{1}(M;V_1,V_2)$ be elliptic on an open
neighbourhood $U$ of $M_+$ and of Dirac type on $M_+$ (for some metric and Hermitian
structures on $M$). Let $Q\in\Psido{-1}(M;V_2,V_1)$ satisfy $\chi(DQ-\Id)\in
\Psido{-\infty}$ for some $\chi\in\CI(M)$ supported in $U$ and equal to $1$ near
$M_+$. Then for $w\in\CI(H;V_2)$:
\begin{enumerate}[label=(\roman*)]
\item $Q(w\,\delta(x))$ restricted to $M_+^\circ$ extends to an element of
$\CI(M_+;V_1)$;
\item the boundary value map $Q_+w=b_+\big(Q(w\,\delta(x))\big)$ is a classical
pseudodifferential operator $Q_+\in\Psido{0}(H;V_2,V_1)$;
\item with $\nu=dx/|dx|$ and $T^*H$ identified with the covectors orthogonal to $\nu$,
\begin{equation*}
\sigma_0(Q_+)(y,\eta)=\frac{i|\eta|\,\sigma(\nu)^*+\sigma(\eta)^*}{2|dx|\,|\eta|}.
\end{equation*}
\end{enumerate}
\end{theorem}
\begin{proof}
Since $\singsupp Q(w\,\delta(x))\subset\supp w\subset H$, only points of $H$ need
attention, and by a partition of unity we may assume $w$ is supported in a coordinate
patch. Choose coordinates $(x,y)$ on an open set $W=(-\epsilon,\epsilon)\times O\subset
\{\chi=1\}$ containing the patch, with $x$ the given function, trivializing $V_1,V_2$
(away from $H$ the analysis is not needed, by pseudolocality). Let $\phi,\psi\in\CIc(W)$ with $\psi=1$ on
$\supp\phi$ and near $\supp w\times\{0\}$, and $\phi=1$ on a neighbourhood of a given
compact $L_0\subset W$. By \cref{ch:manifolds}, in these coordinates $\phi Q\psi=
\Op(q)+R$ with $q\in S^{-1}(\R^n;\R^n)$ and $R$ with smooth compactly supported kernel,
and $\phi Q\psi(w\delta(x))=\phi Q(w\delta(x))$. Moreover, by
\eqref{boundary:eq:comp-diff}, $\Op(d\circ q)=D\phi Q\psi-DR=\phi\psi+\phi(DQ-\Id)\psi
+[D,\phi]Q\psi-DR$; the commutator is multiplication by a matrix function vanishing
on $\{\phi=1\}^\circ$, so over the compact set $K'=L'\cap\{x\ge0\}$, for
$L'\subset\{\phi=1\}^\circ$ a compact neighbourhood of $L_0$, we have $d\circ q-1\in
S^{-\infty}(K')$. \Cref{boundary:lem:rational} applies with $K=\{x\ge0\}$ and gives
$q=\rho_N+e$ with $e\in S^{-N-1}(L)$, $L=L_0\cap\{x\ge0\}$.

Let $\phi_1\in\CIc(W)$ equal $1$ near $L_0$. For $z\in L$, $\Op(q)f(z)=
\Op(\phi_1\rho_N)f(z)+\Op(e')f(z)$ with $e'=q-\phi_1\rho_N\in S^{-1}$, $e'=e$ near
$L$. We compute both terms on $f=w\otimes\delta(x)$, $\hat f(\xi,\eta)=\hat w(\eta)$,
using the approximations $f_\epsilon=w\otimes\varphi_\epsilon(x)$,
$\varphi_\epsilon=\epsilon^{-1}\varphi(x/\epsilon)$, which converge to $f$ in
$\cSp$, so that $\Op(a)f_\epsilon\to\Op(a)f$ in $\cSp$ for any symbol $a$.

\emph{The remainder.} For $z\in L$, $|D_z^\alpha(e^{iz\cdot\zeta}e'(z,\zeta))|\le
C\ang\zeta^{|\alpha|-N-1}$ and $\int\ang{\zeta}^{|\alpha|-N-1}d\xi\le C\ang\eta^{|\alpha|
-N}$ if $|\alpha|\le N-1$. Since $\hat w$ is rapidly decreasing, dominated convergence
shows that $\Op(e')f_\epsilon(z)\to\int\!\!\int e^{iz\cdot\zeta}e'(z,\zeta)\hat
w(\eta)\,\dbar\xi\,\dbar\eta$ uniformly on $L$, together with $z$-derivatives of order
$\le N-1$. Thus on the interior of $L$ (relative to $\R^n$), $\Op(e')f$ is this
absolutely convergent integral, which is $\mathcal{C}^{N-1}$ on $L$, including up to
$x=0$. Its value at $x=0$ is $\int e^{iy\cdot\eta}t_N(y,\eta)\hat w(\eta)\dbar\eta$
with $t_N(y,\eta)=\int e((0,y),\xi,\eta)\dbar\xi$, and since $D^\beta_\eta e\in
S^{-N-1-|\beta|}(L)$, $t_N\in S^{-N}$.

\emph{The main term.} $\phi_1\rho_N$ is a finite sum of terms $\phi_1\rho$, $\rho\in
\Rat{-1}$. By \cref{boundary:lem:contour}(i), $\int e^{ix\xi}\rho(z,\xi,\eta)\hat
\varphi(\epsilon\xi)\dbar\xi$ is the convolution in $x$ of
$H(x)G^+_\rho-H(-x)G^-_\rho$ with $\varphi_\epsilon$; for $x>0$ and $\epsilon$ small
it converges to $G^+_\rho(z,x,\eta)$, with error at most $C\epsilon\ang\eta^{r+2}$
uniformly for $x\ge2\epsilon$ by (ii) of that lemma. Hence for $z\in L$ with $x>0$,
\begin{equation*}
\Op(\phi_1\rho)f(z)=\phi_1(z)\int e^{iy\cdot\eta}G^+_\rho(z,x,\eta)\hat w(\eta)\,
\dbar\eta,
\end{equation*}
and by \cref{boundary:lem:contour}(ii) and the rapid decrease of $\hat w$ the right side
is smooth in $z\in L$ up to $x=0$, with boundary value $\int e^{iy\cdot\eta}
G^+_\rho((0,y),0,\eta)\hat w(\eta)\dbar\eta$.

Combining, $Q(w\delta(x))$ is $\mathcal{C}^{N-1}$ up to $x=0$ near $L$ for every $N$,
which proves (i) since $L_0$ was arbitrary. Moreover, near $L$,
\begin{equation*}
Q_+w=\Op_H\big(\sigma^{(N)}+t_N\big)w+R_+w,\qquad\sigma^{(N)}(y,\eta)=
\sum G^+_{\phi_1\rho}((0,y),0,\eta),
\end{equation*}
with $R_+$ smoothing, the sum over the terms of $\rho_N$. Here $\sigma^{(N)}\in S^0$
is classical, $t_N\in S^{-N}$, and $\sigma^{(N+1)}-\sigma^{(N)}\in S^{-N}$ since
$\rho_{N+1}-\rho_N\in\Rat{-N-1}$. Hence locally $Q_+$ is a pseudodifferential operator
of order $0$ whose symbol agrees modulo $S^{-N}$ with a classical symbol for every
$N$, i.e.\ it is classical. If $\psi_1,\phi_2\in\CI(H)$ have disjoint supports,
$\psi_1Q_+\phi_2$ has smooth kernel, since $Q(\phi_2w\delta(x))$ is given near
$\supp\psi_1$ by integrating the kernel of $Q$, smooth off the diagonal, against
$\phi_2w$ over $H$. Choosing a cover of $H$ by coordinate patches such that any two
that meet lie in a common patch, it follows that $Q_+\in\Psido{0}(H;V_2,V_1)$, which
is (ii).

For (iii) we compute at a point of $H$. Changing the coordinates $y$ (which does not
affect $x$, $\delta(x)$ or $Q_+$) we may assume $dy_k\perp dx$ at that point, i.e.\
$\beta=0$; then $g^{11}=|dx|^2$, $g'(\eta)=|\eta|^2$ and $\xi^0_\pm=\pm i|\eta|/|dx|$.
The principal symbol of $Q_+$ comes from $\rho_1=\tau/(1+g)$ only, and by
\cref{boundary:lem:contour}(iv) equals
$i\tau(\xi^0_+dx+\eta)/(2i|dx||\eta|)$; since $\xi^0_+dx=i|\eta|\nu$ and $\tau$ is
complex linear, this is the stated formula.
\end{proof}

\begin{remark}\label{boundary:rem:other-side}
If $D$ is of Dirac type on $M_-=\{x\le0\}$ as well, the same argument gives
smoothness of $Q(w\delta(x))$ up to $H$ from $M_-$, with boundary value
$Q_-\in\Psido{0}(H;V_2,V_1)$, $\sigma_0(Q_-)=(-i|\eta|\sigma(\nu)^*+\sigma(\eta)^*)/
(2|dx||\eta|)$, obtained from the residue at $\xi_-$.
\end{remark}

\section{The Calder\'on projector for a dividing hypersurface}\label{boundary:sec:calderon-dividing}

We return to the standing hypotheses of \cref{boundary:sec:dividing}. By
\cref{boundary:thm:layer} and \cref{boundary:rem:other-side} with $Q=D^{-1}$, for
$w\in\CI(H;V_2)$ the distribution $D^{-1}(w\delta(x))$ is smooth up to $H$ from both
sides, which completes the proof of \cref{boundary:prop:layer-iso}, and the boundary
values
\begin{equation*}
Q_\pm w=\lim_{\pm x\downarrow0}D^{-1}(w\,\delta(x))
\end{equation*}
define $Q_\pm\in\Psido{0}(H;V_2,V_1)$. Applying \eqref{boundary:eq:jump-solutions} to
$u=D^{-1}(w\delta(x))\in\mathcal{N}_H$ gives $\frac1i\sigma(dx)(Q_+-Q_-)w=w$, i.e.
\begin{equation}\label{boundary:eq:Qjump}
Q_+-Q_-=i\,\sigma(dx)^{-1}.
\end{equation}

\begin{proof}[Proof of \cref{boundary:thm:calderon-dividing}]
Define
\begin{equation}\label{boundary:eq:PiC-def}
\PiC v=Q_+\Big(\frac1i\sigma(dx)v\Big)=b_+D^{-1}\Big(\frac1i\sigma(dx)v\,\delta(x)\Big),
\quad v\in\CI(H;V),
\end{equation}
so $\PiC\in\Psido{0}(H;V)$ is classical, and by \eqref{boundary:eq:Qjump}
$\Id-\PiC=-Q_-\frac1i\sigma(dx)$.

If $u\in\mathcal{N}_+$ then $u_c\in\mathcal{N}_H$ (zero on $M_-$) and by
\eqref{boundary:eq:jump-solutions} $Du_c=\frac1i\sigma(dx)(b_+u)\delta(x)$, so
$u_c=D^{-1}(\frac1i\sigma(dx)(b_+u)\delta(x))$ and, taking boundary values from
$M_+$,
\begin{equation}\label{boundary:eq:PiC-on-data}
\PiC\,b_+u=b_+u,\quad u\in\mathcal{N}_+.
\end{equation}
Conversely for any $v\in\CI(H;V)$, $u=D^{-1}(\frac1i\sigma(dx)v\delta(x))|_{M_+}$
lies in $\mathcal{N}_+$ and $b_+u=\PiC v$. Hence $\Ran\PiC=b_+(\mathcal{N}_+)$, and
by \eqref{boundary:eq:PiC-on-data} $\PiC^2v=\PiC v$. Together with
\cref{boundary:lem:unique}, $b_+:\mathcal{N}_+\to\PiC\CI(H;V)$ is an isomorphism.

Similarly, for $u\in\mathcal{N}_-$, $Du_c=-\frac1i\sigma(dx)(b_-u)\delta(x)$, so
$b_-u=-Q_-\frac1i\sigma(dx)b_-u=(\Id-\PiC)b_-u$; and for any $v$, $-D^{-1}(\frac1i
\sigma(dx)v\delta(x))|_{M_-}\in\mathcal{N}_-$ has boundary value $(\Id-\PiC)v$. Hence
$b_-:\mathcal{N}_-\to(\Id-\PiC)\CI(H;V)$ is an isomorphism.

Uniqueness: a projection on $\CI(H;V)$ is determined by its range and null space,
and these are forced to be $b_+(\mathcal{N}_+)$ and $b_-(\mathcal{N}_-)$; a
pseudodifferential operator is determined by its action on $\CI$.

Finally, by \cref{boundary:thm:layer}(iii), $\sigma_0(\PiC)=\sigma_0(Q_+)\frac1i|dx|
\sigma(\nu)=\frac12\big(\sigma(\nu)^*\sigma(\nu)+\sigma(\eta)^*\sigma(\nu)/(i|\eta|)
\big)$. By \eqref{boundary:eq:dirac} $\sigma(\nu)^*\sigma(\nu)=\Id$, and by
\eqref{boundary:eq:clifford}, since $g(\eta,\nu)=0$,
$\sigma(\eta)^*\sigma(\nu)=-\sigma(\nu)^*\sigma(\eta)=-\sigma(\nu)^{-1}\sigma(\eta)$,
which gives \eqref{boundary:eq:calderon-symbol}. With $A=i\sigma(\nu)^{-1}\sigma(\eta)/
|\eta|=i\sigma(\nu)^*\sigma(\eta)/|\eta|$, the second Clifford relation gives
$\sigma(\eta)\sigma(\nu)^*=-\sigma(\nu)\sigma(\eta)^*$, so $A^2=-\sigma(\nu)^*\sigma(\eta)
\sigma(\nu)^*\sigma(\eta)/|\eta|^2=\sigma(\nu)^*\sigma(\nu)\sigma(\eta)^*\sigma(\eta)/
|\eta|^2=\Id$, and $A^*=-i\sigma(\eta)^*\sigma(\nu)/|\eta|=A$.
\end{proof}

\begin{example}
For $D=D_x+iD_y$ (\cref{boundary:ex:dirac}) we have $\sigma(\nu)=1$, $\sigma(\eta)=
i\eta$, so $\sigma_0(\PiC)=\frac12(1-\eta/|\eta|)$, the projection onto negative
frequencies. This is as it should be: $Du=0$ means $u$ is holomorphic in $x+iy$, and
$e^{\eta(x+iy)}$ is bounded in $x>0$ precisely when $\eta<0$.
\end{example}

\subsection*{The Poisson operator and a right inverse}
The proof shows that
\begin{equation*}
I_D:\PiC\CI(H;V)\to\mathcal{N}_+,\quad I_Dv=D^{-1}\Big(\frac1i\sigma(dx)v\,\delta(x)
\Big)\Big|_{M_+},
\end{equation*}
is the inverse of $b_+$: $u=I_Dv$ is the unique solution of $Du=0$ in $M_+^\circ$,
$u\in\CI(M_+;V_1)$, $b_+u=v$. This is the \emph{Poisson operator} for the Calder\'on
boundary condition. As well as $I_D$ we have a natural right inverse for $D$ on $M_+$.

\begin{lemma}\label{boundary:lem:RD}
If $f\in\CI(M_+;V_2)$ then $R_Df=D^{-1}(f_c)|_{M_+}\in\CI(M_+;V_1)$, and $D\circ
R_D=\Id$ on $\CI(M_+;V_2)$.
\end{lemma}
\begin{proof}
Certainly $D(D^{-1}f_c)=f_c$, so $DR_Df=f$ in $M_+^\circ$; the point is the
regularity. By \cref{boundary:lem:formal}(i) there is $u'\in\CI(M_+;V_1)$ with
$b_+u'=0$ and $h=Du'-f\in\dCI(M_+;V_2)$. By the jump formula $D(u'_c)=(Du')_c=f_c+h_c$,
where $h_c\in\CI(M;V_2)$. Hence $D^{-1}f_c=u'_c-D^{-1}h_c$ and $R_Df=u'-(D^{-1}h_c)
|_{M_+}\in\CI(M_+;V_1)$.
\end{proof}

$R_D$ is not a two-sided inverse: $D$ has the large null space $\mathcal{N}_+=I_D\big(
\PiC\CI(H;V)\big)$ on $M_+$. Combining the two operators solves the boundary problem
with Calder\'on boundary condition, see \cref{boundary:ex:calderon-bvp}.

\section{Smoothing operators and generalized inverses}\label{boundary:sec:smoothing}

On a manifold with boundary there is no invertible extension of $D$ at hand, and we
must work with parametrices; to correct their errors we use some abstract functional
analysis. Let $\mathcal H$ be a separable Hilbert space, $\Lin=\Lin(\mathcal H)$,
$\Kpt=\Kpt(\mathcal H)$ the compact operators and $\Fin=\Fin(\mathcal H)$ the
finite-rank operators, so $\Kpt$ is the norm closure of $\Fin$.

\begin{definition}\label{boundary:def:biideal}
A subspace $\cJ\subset\Lin$ with a Fr\'echet topology given by norms
$\norm{\cdot}_k$, $k\in\N_0$, is a \emph{bi-ideal} if:
\begin{enumerate}[label=(\alph*)]
\item\label{boundary:it:alg} $\cJ$ is a subalgebra and for each $k$ there are $k'$ and
$C_k$ with $\norm{AB}_k\le C_k\norm{A}_{k'}\norm{B}_{k'}$, $A,B\in\cJ$;
\item\label{boundary:it:star} $\norm{A}_{\Lin}\le C\norm{A}_0$ and $A^*\in\cJ$ with
$\norm{A^*}_k\le C_k\norm{A}_{k'}$ for $A\in\cJ$;
\item\label{boundary:it:dense} $\Fin\cap\cJ$ is dense in $\cJ$;
\item\label{boundary:it:corner} if $A_1,A_2\in\cJ$ and $B\in\Lin$ then $A_1BA_2\in\cJ$,
with $\norm{A_1BA_2}_k\le C_k\norm{A_1}_{k'}\norm{B}_\Lin\norm{A_2}_{k'}$.
\end{enumerate}
\end{definition}

By \ref{boundary:it:star} and \ref{boundary:it:dense}, $\cJ\subset\Kpt$. Note that
$\cJ$ need not be an ideal in $\Lin$; \ref{boundary:it:corner} replaces this.

\begin{proposition}\label{boundary:prop:smoothing-biideal}
Let $X$ be a compact manifold, with or without boundary, with a smooth positive
density, and $E$ a Hermitian bundle. The operators on $\mathcal{H}=L^2(X;E)$ with
kernels in $\CI(X\times X;\Hom(E))$, $Au(z)=\int_XA(z,z')u(z')$, with the norms
$\norm{A}_k=\norm{A(\cdot,\cdot)}_{\mathcal{C}^k(X\times X)}$ (defined using a fixed
finite atlas), form a bi-ideal.
\end{proposition}
\begin{proof}
The kernel of $AB$ is $\int_XA(z,z'')B(z'',z')dz''$, whose $\mathcal{C}^k$ norm is
bounded by $\operatorname{vol}(X)\norm{A}_k\norm{B}_k$; the kernel of $A^*$ is
$A(z',z)^*$; and $\norm{A}_\Lin\le\norm{A(\cdot,\cdot)}_{L^2(X\times X)}\le C\norm{A}_0$.
For \ref{boundary:it:corner}, the kernel of $A_1BA_2$ is $\inner{BA_2(\cdot,z')}{A_1(z,
\cdot)^*}_{L^2}$, and derivatives in $z,z'$ fall on $A_1$, $A_2$, giving the estimate.
For \ref{boundary:it:dense}, embed $X$ in its double, extend the kernel to a smooth
kernel on $\Xt\times\Xt$, and by a partition of unity reduce to a kernel in
$\CIc(\R^n\times\R^n)$; regarding it as a smooth function on a large torus, its Fourier
partial sums $\sum c_{jk}e^{ij\cdot z}e^{ik\cdot z'}$ converge in $\CI$, and after
multiplying by $\chi(z)\chi(z')$, $\chi=1$ on the support, they give finite sums of
products $\varphi(z)\psi(z')$, i.e.\ finite-rank operators in $\cJ$.
\end{proof}

The elements of this bi-ideal are the \emph{smoothing operators}; they have the
following characterization (in the boundaryless case $\dCmI=\CmI$).

\begin{proposition}\label{boundary:prop:smoothing-char}
A linear operator $A:\dCI(X;E)\to\CmI(X;F)$ has kernel in $\CI(X\times X;\Hom(E,F)
\otimes\Omega_R)$ if and only if it extends to a continuous map $\dCmI(X;E)\to
\CI(X;F)$, i.e.\ to continuous maps $\Hdot{-m}(X;E)\to\CI(X;F)$ for all $m$.
\end{proposition}
\begin{proof}
If $A$ has smooth kernel, $Au(z)=u(A(z,\cdot))$ makes sense for supported
distributions $u$ (the test functions being smooth up to the boundary) and is smooth
and continuous in $u$. Conversely, for $z'\in X$ and $e\in E_{z'}$, $\delta_{z'}e\in
\Hdot{-m}(X;E)$ when $m>n/2$, and $z'\mapsto\delta_{z'}$ is $\mathcal{C}^k$ into
$\Hdot{-m}$ for $m>n/2+k$ (as seen from the Fourier transform $e^{-iz'\cdot\zeta}$ in
local coordinates). Thus $A(z,z')e=(A\delta_{z'}e)(z)$ is $\mathcal{C}^k$ in $z'$
with values in $\CI(X;F)$ for every $k$, hence smooth on $X\times X$. For
$u\in\dCI(X;E)$, $u=\int u(z')\delta_{z'}dz'$ as a Riemann integral in $\Hdot{-m}$,
so by continuity $Au=\int A(\cdot,z')u(z')dz'$.
\end{proof}

\begin{proposition}\label{boundary:prop:gen-inverse}
If $\cJ$ is a bi-ideal and $A\in\cJ$, then $\Id+A$ has a generalized inverse $\Id+B$,
$B\in\cJ$:
\begin{equation*}
(\Id+B)(\Id+A)=\Id-\pi_L,\quad(\Id+A)(\Id+B)=\Id-\pi_R,
\end{equation*}
where $\pi_L,\pi_R\in\cJ\cap\Fin$ are the orthogonal projections onto $\Nul(\Id+A)$ and
onto $\Ran(\Id+A)^\perp=\Nul(\Id+A^*)$.
\end{proposition}
\begin{proof}
As $A$ is compact, $\Id+A$ is Fredholm. If $u\in\Nul(\Id+A)$ then $u=-Au$, so
$\pi_L=-A\pi_L$, and taking adjoints $\pi_L=-\pi_LA^*$; hence $\pi_L=A\pi_LA^*\in\cJ$
by \ref{boundary:it:corner}. Similarly $\pi_R=A^*\pi_RA\in\cJ$. Let $G\in\Lin$ be
the inverse of $\Id+A:\Nul(\Id+A)^\perp\to\Ran(\Id+A)$, extended by $0$ on
$\Ran(\Id+A)^\perp$; then $G(\Id+A)=\Id-\pi_L$ and $(\Id+A)G=\Id-\pi_R$. Setting
$B=G-\Id$, these identities give $B=-\pi_L-A-BA$ and $B=-\pi_R-A-AB$; substituting
the second into the first,
\begin{equation*}
B=-\pi_L-A+\pi_RA+A^2+ABA,
\end{equation*}
and each term lies in $\cJ$ by \ref{boundary:it:alg} and \ref{boundary:it:corner}.
\end{proof}

Now let $\mathcal{H}_1,\mathcal{H}_2$ be Hilbert spaces, $\cJ_i\subset\Lin(\mathcal
H_i)$ bi-ideals and $A\in\Lin(\mathcal H_1,\mathcal H_2)$. A \emph{left parametrix}
for $A$ modulo $\cJ_1$ is $B_L\in\Lin(\mathcal H_2,\mathcal H_1)$ with $B_LA=\Id+J_L$,
$J_L\in\cJ_1$; a \emph{right parametrix} modulo $\cJ_2$ is $B_R$ with $AB_R=\Id+J_R$,
$J_R\in\cJ_2$.

\begin{proposition}\label{boundary:prop:left-param}
If $A$ has a left parametrix $B_L$ modulo $\cJ_1$, then $A$ has closed range and
finite-dimensional null space, the orthogonal projection $\pi_L$ onto $\Nul A$ lies in
$\cJ_1\cap\Fin$, and there are $J\in\cJ_1$, $\pi'_L\in\cJ_1\cap\Fin$ a self-adjoint
projection with $\pi_L\le\pi'_L$, and $G\in\Lin(\mathcal{H}_2,\mathcal{H}_1)$ such that
\begin{equation*}
B'_L=(\Id+J)B_L\ \text{satisfies}\ B'_LA=\Id-\pi'_L,\quad\text{and}\quad
B''_L=B'_L+(\pi'_L-\pi_L)G\ \text{satisfies}\ B''_LA=\Id-\pi_L,
\end{equation*}
with $\pi'_L-\pi_L\in\cJ_1\cap\Fin$.
\end{proposition}
\begin{proof}
By \cref{boundary:prop:gen-inverse} there is $J\in\cJ_1$ with $(\Id+J)(\Id+J_L)=
\Id-\pi'_L$, which gives $B'_L$. Then $\Nul A\subset\Ran\pi'_L$ is finite dimensional
and $\pi_L=\pi'_L\pi_L\pi'_L\in\cJ_1$. On the closed subspace $\Ran(\Id-\pi'_L)$,
$u=B'_LAu$, so $\norm{u}\le\norm{B'_L}\norm{Au}$ and $A(\Ran(\Id-\pi'_L))$ is closed;
adding the finite-dimensional space $A(\Ran\pi'_L)$, $\Ran A$ is closed. Let $G$ be
the generalized inverse of $A$ ($GA=\Id-\pi_L$, $G=0$ on $(\Ran A)^\perp$). Since
$\pi'_L\pi_L=\pi_L$, $B''_LA=\Id-\pi'_L+(\pi'_L-\pi_L)(\Id-\pi_L)=\Id-\pi_L$, and
$\pi'_L-\pi_L=\pi'_L(\Id-\pi_L)\pi'_L\in\cJ_1$.
\end{proof}

\begin{proposition}\label{boundary:prop:right-param}
If $A$ has a right parametrix $B_R$ modulo $\cJ_2$, then $A$ has closed range of finite
codimension, the orthogonal projection $\pi_R$ onto $(\Ran A)^\perp$ lies in
$\cJ_2\cap\Fin$, and there are $J\in\cJ_2$ and a self-adjoint projection $\pi'_R\in
\cJ_2\cap\Fin$ with $\pi_R\le\pi'_R$ such that $B'_R=B_R(\Id+J)$ satisfies
$AB'_R=\Id-\pi'_R$; moreover $B''_R=B'_R+G(\pi'_R-\pi_R)$, $G$ the generalized inverse
of $A$, satisfies $AB''_R=\Id-\pi_R$.
\end{proposition}
\begin{proof}
Choose $J$ by \cref{boundary:prop:gen-inverse} with $(\Id+J_R)(\Id+J)=\Id-\pi'_R$.
Then $\Ran A\supset\Ran(\Id-\pi'_R)$, a closed subspace of finite codimension, so
$\Ran A$ is closed of finite codimension, $(\Ran A)^\perp\subset\Ran\pi'_R$ and
$\pi_R=\pi'_R\pi_R\pi'_R\in\cJ_2$. Finally $AG=\Id-\pi_R$ and $\pi_R\pi'_R=\pi_R$ give
$AB''_R=\Id-\pi'_R+(\Id-\pi_R)(\pi'_R-\pi_R)=\Id-\pi_R$.
\end{proof}

An operator with closed range and either finite-dimensional null space or
finite-codimensional range is called \emph{semi-Fredholm}. The existence of a left or
right parametrix modulo compact operators is necessary and sufficient for these two
kinds of semi-Fredholmness.

\section{Manifolds with boundary: a right inverse}\label{boundary:sec:right-inverse}

\noindent\textbf{Standing hypotheses (boundary case).} From now on $X$ is a compact
Riemannian manifold with boundary, $V_1,V_2$ are Hermitian bundles over $X$, and
$D\in\Diff{1}(X;V_1,V_2)$ is a generalized Dirac operator. We embed $X$ in its double
$\Xt$, a compact manifold without boundary, with $x\in\CI(\Xt)$ a function such that
$X=\{x\ge0\}$ and $dx\ne0$ on $\pa X=\{x=0\}$, and extend the metric, the Hermitian
bundles and $D$ (by extending coefficients across $\pa X$ in a collar and cutting
off) to $\tilde D\in\Diff{1}(\Xt;\tilde V_1,\tilde V_2)$. Since ellipticity is an open
condition, $\tilde D$ is elliptic on an open neighbourhood $U$ of $X$; it is of Dirac
type on $X$, but need not be elsewhere. By the parametrix construction of
\cref{ch:manifolds} (which is local), there is $\tilde Q\in\Psido{-1}(\Xt;\tilde V_2,
\tilde V_1)$ and $\chi\in\CI(\Xt)$ with $\supp\chi\subset U$, $\chi=1$ near $X$, such
that
\begin{equation*}
E_R=\chi(\tilde D\tilde Q-\Id),\quad E_L=(\tilde Q\tilde D-\Id)\chi\quad\text{are
smoothing.}
\end{equation*}
We also use a smooth positive density on $X$ (the Riemannian one) and the
corresponding $L^2$ spaces. Finally we assume \emph{unique continuation}:
\begin{equation}\tag{UC}\label{boundary:eq:UC}
\text{if}\ u\in\dCI(X;V_1),\ Du=0\ \text{then}\ u=0;\quad\text{if}\ f\in\dCI(X;V_2),\
D^*f=0\ \text{then}\ f=0.
\end{equation}

\begin{remark}
\eqref{boundary:eq:UC} fails if $X$ has a component without boundary on which $D$
or $D^*$ has null space. When every component of $X$ meets $\pa X$ it holds for
operators of Dirac type: since $D^*D$ and $DD^*$ have scalar principal symbol
$|\zeta|^2$, this follows from the unique continuation theorem of Aronszajn for such
second-order operators, applied after extending the Dirac structure a little across
the boundary near a boundary point; see the Notes. We simply assume it.
\end{remark}

For $f\in L^2(X;V_2)$ set
\begin{equation*}
Q'f=R_X\tilde Q(f_c)\in H^1(X;V_1),
\end{equation*}
where $f_c\in L^2(\Xt;\tilde V_2)$ is the extension by zero; $Q':L^2(X;V_2)\to
H^1(X;V_1)$ is bounded.

\begin{lemma}\label{boundary:lem:Qprime}
$Q'$ maps $\CI(X;V_2)$ into $\CI(X;V_1)$, and $DQ'=\Id+E'$ where $E'f=R_XE_R(f_c)$
has kernel in $\CI(X\times X;\Hom(V_2))$.
\end{lemma}
\begin{proof}
The second statement holds since $\chi=1$ on $X$: $DQ'f=R_X\tilde D\tilde Qf_c=
f+R_XE_Rf_c$. For the first, as in \cref{boundary:lem:RD}, choose by
\cref{boundary:lem:formal} $u'\in\CI(X;V_1)$ with $bu'=0$ and $h=Du'-f\in\dCI(X;V_2)$.
By the jump formula $\tilde D(u'_c)=f_c+h_c$ with $h_c$ smooth, so
$u'_c+E_Lu'_c=\tilde Q\tilde D(u'_c)=\tilde Qf_c+\tilde Qh_c$, and $Q'f=u'+R_X(E_Lu'_c
-\tilde Qh_c)\in\CI(X;V_1)$.
\end{proof}

By \cref{boundary:prop:smoothing-biideal}, $E'$ lies in the bi-ideal of smoothing
operators on $L^2(X;V_2)$, so $Q'$ is a right parametrix for $D:H^1(X;V_1)\to
L^2(X;V_2)$ and \cref{boundary:prop:right-param} shows that this map has closed range
of finite codimension.

\begin{proposition}\label{boundary:prop:surjective}
$D:\CI(X;V_1)\to\CI(X;V_2)$ is surjective if and only if there is no non-trivial
$f\in\dCI(X;V_2)$ with $D^*f=0$. In that case $D:H^1(X;V_1)\to L^2(X;V_2)$ is
surjective as well.
\end{proposition}
\begin{proof}
If $f\in\dCI(X;V_2)$ and $D^*f=0$ then, integrating by parts (no boundary term
arises as $f$ vanishes at $\pa X$), $\inner{Du}{f}=\inner{u}{D^*f}=0$ for all $u\in
\CI(X;V_1)$, so $D$ is not surjective unless $f=0$.

Conversely, assume there is no such $f$, and let $f\in L^2(X;V_2)$ be orthogonal to
$D(H^1(X;V_1))$. For $\varphi\in\CIc(U;\tilde V_1)$, $\inner{\tilde D\varphi}{f_c}=
\inner{D(\varphi|_X)}{f}=0$, so $\tilde D^*f_c=0$ on $U$. By elliptic regularity $f_c$
is smooth on $U$; as it vanishes outside $X$, $f\in\dCI(X;V_2)$ and $D^*f=0$, so
$f=0$. Thus $D:H^1\to L^2$ is surjective.

For the smooth statement, apply \cref{boundary:prop:right-param}: there are a
smoothing $J$ and a finite-rank smoothing self-adjoint projection $\pi'_R$, with
range in $\CI(X;V_2)$, such that $DQ''=\Id-\pi'_R$, $Q''=Q'(\Id+J)$. Since $J$ maps
$L^2$ into $\CI$, $Q''$ maps $\CI(X;V_2)$ into $\CI(X;V_1)$, so the image $W$ of
$\CI(X;V_1)$ under $D$ contains $(\Id-\pi'_R)\CI(X;V_2)$. Hence $f\in W$ if and only if
$\pi'_Rf\in W\cap\Ran\pi'_R$. If $W\ne\CI(X;V_2)$, choose a non-zero linear functional
$\lambda$ on the finite-dimensional space $\Ran\pi'_R$ vanishing on $W\cap\Ran\pi'_R$;
then $\ell=\lambda\circ\pi'_R$ is $L^2$-continuous, vanishes on $W$, but not on
$\Ran\pi'_R\subset\CI$. But $W$ is dense in $D(H^1)=L^2$, since $\CI$ is dense in
$H^1$, so $\ell=0$, a contradiction.
\end{proof}

\begin{theorem}\label{boundary:thm:right-inverse}
Under the standing hypotheses (including \eqref{boundary:eq:UC}) $D$ has a right inverse
\begin{equation*}
Q:\CI(X;V_2)\to\CI(X;V_1),\quad DQ=\Id,\quad Q=Q'+E'',
\end{equation*}
with $E''$ smoothing; $Q$ extends to a bounded map $L^2(X;V_2)\to H^1(X;V_1)$ and is a
right inverse there too.
\end{theorem}
\begin{proof}
With $Q''$ and $\pi'_R=\sum_je_j\inner{\cdot}{e_j}$ ($e_j\in\CI$ orthonormal) as in the
previous proof, choose by \cref{boundary:prop:surjective} $v_j\in\CI(X;V_1)$ with
$Dv_j=e_j$, and set $Q=Q''+\sum_jv_j\inner{\cdot}{e_j}$. Then $DQ=\Id-\pi'_R+\pi'_R=
\Id$. $Q-Q'=Q'J+\sum_jv_j\inner\cdot{e_j}$ has smooth kernel, since $Q'$ is continuous
on $\CI(X)$ (closed graph theorem) and $J$ has smooth kernel.
\end{proof}

\section{The boundary map on solutions}\label{boundary:sec:boundary-map}

For $s\in\R$ let
\begin{equation*}
\mathcal{N}^s=\{u\in H^s(X;V_1);\ Du=0\ \text{in}\ X^\circ\},\qquad
\mathcal{N}^\infty=\{u\in\CI(X;V_1);\ Du=0\}=\bigcap_s\mathcal{N}^s.
\end{equation*}
By \cref{boundary:lem:formal}(ii) and \eqref{boundary:eq:UC}, $b:\mathcal{N}^\infty\to
\CI(\pa X;V_1)$ is injective; its range
\begin{equation*}
\Cau=b(\mathcal{N}^\infty)\subset\CI(\pa X;V_1)
\end{equation*}
is the space of \emph{Cauchy data}. To show that it is the range of a
pseudodifferential projection we extend $b$ to $L^2$ solutions. Let
\begin{equation*}
\mathcal{D}=\{u\in L^2(X;V_1);\ Du\in L^2(X;V_2)\}
\end{equation*}
with the graph norm, where $Du$ is computed in $\CmI(X;V_2)$.

\begin{lemma}\label{boundary:lem:beta}
For $u\in\mathcal{D}$ there is a unique $\beta(u)\in H^{-\frac12}(\pa X;V_1)$ such that
\begin{equation}\label{boundary:eq:beta}
\tilde D(u_c)=(Du)_c+\frac1i\sigma(dx)\beta(u)\,\delta(x),
\end{equation}
and $\beta:\mathcal{D}\to H^{-\frac12}(\pa X;V_1)$ is continuous with $\beta(u)=bu$ for
$u\in\CI(X;V_1)$. In particular $\tilde b=\beta|_{\mathcal{N}^0}:\mathcal{N}^0\to
H^{-\frac12}(\pa X;V_1)$ is a continuous injective extension of $b$.
\end{lemma}
\begin{proof}
$T=\tilde D(u_c)-(Du)_c\in H^{-1}(\Xt;\tilde V_2)$ depends continuously on $u\in
\mathcal{D}$ and is supported in $\pa X$: it vanishes outside $X$ and equals $Du-Du=0$
in $X^\circ$. By \cref{boundary:lem:layers}, $T=w\,\delta(x)$ with $w\in
H^{-\frac12}$ depending continuously on $T$; set $\beta(u)=i\sigma(dx)^{-1}w$. For
smooth $u$ this is $bu$ by \cref{boundary:cor:jump1}. If $u\in\mathcal{N}^0$ and
$\tilde b u=0$, then $\tilde D(u_c)=0$ on $U$, so $u_c$ is smooth by elliptic
regularity, whence $u\in\dCI(X;V_1)$, $Du=0$, and $u=0$ by \eqref{boundary:eq:UC}.
\end{proof}

For $v\in H^{-\frac12}(\pa X;V_1)$ the layer $T_v=\frac1i\sigma(dx)v\,\delta(x)$ lies
in $H^{-1}(\Xt;\tilde V_2)$, and we set
\begin{equation*}
\tilde I_Dv=R_X\tilde Q(T_v)\in L^2(X;V_1).
\end{equation*}
Since $R_XT_v=0$ and $\chi=1$ on $X$,
\begin{equation}\label{boundary:eq:Rprime}
D\tilde I_Dv=R_X\tilde D\tilde QT_v=R_XE_RT_v=:R'v,
\end{equation}
where $R'$ has smooth kernel on $X\times\pa X$, and so extends to a continuous map
$\CmI(\pa X;V_1)\to\CI(X;V_2)$. In particular $\tilde I_D:H^{-\frac12}(\pa X;V_1)\to
\mathcal{D}$ is continuous. Applying \cref{boundary:thm:layer} to $\tilde D$ on $\Xt$,
with $M_+=X$ and $Q=\tilde Q$, we obtain the classical operator
\begin{equation}\label{boundary:eq:P}
P=Q_+\circ\frac1i\sigma(dx)\in\Psido{0}(\pa X;V_1),\quad Pv=b\,\tilde I_Dv\ \ (v\in
\CI(\pa X;V_1)),
\end{equation}
with $\tilde I_Dv\in\CI(X;V_1)$ for smooth $v$, and $\sigma_0(P)$ given by the right
side of \eqref{boundary:eq:calderon-symbol} (the computation in the proof of
\cref{boundary:thm:calderon-dividing} only uses the Dirac identities at $\pa X$).

\begin{lemma}\label{boundary:lem:repr}
\begin{enumerate}[label=(\roman*)]
\item $\beta(\tilde I_Dv)=Pv$ for all $v\in H^{-\frac12}(\pa X;V_1)$.
\item (Representation formula) For $u\in\mathcal{D}$,
\begin{equation*}
u=Q'(Du)+\tilde I_D\beta(u)+Gu,\qquad Gu=-R_XE_L(u_c),
\end{equation*}
where $G$ has smooth kernel on $X\times X$; in particular $u=\tilde I_D\tilde b u+Gu$
for $u\in\mathcal{N}^0$.
\item There is $C$ such that $\norm{u}_{L^2}\le C\norm{\tilde bu}_{H^{-\frac12}}$ for
$u\in\mathcal{N}^0$. Hence $\overline{\Cau}:=\tilde b(\mathcal{N}^0)$ is closed in
$H^{-\frac12}(\pa X;V_1)$ and $\tilde b:\mathcal{N}^0\to\overline{\Cau}$ is an
isomorphism.
\end{enumerate}
\end{lemma}
\begin{proof}
(i) holds for smooth $v$ by \eqref{boundary:eq:P}, and both sides are continuous on
$H^{-\frac12}$. (ii) By \eqref{boundary:eq:beta}, $\tilde D(u_c)=(Du)_c+T_{\beta(u)}$;
applying $\tilde Q$ and using $\tilde Q\tilde D(u_c)=u_c+E_Lu_c$ gives the formula on
restricting to $X$. (iii) If not, there are $u_k\in\mathcal{N}^0$ with $\norm{u_k}=1$
and $\tilde bu_k\to0$. As $G$ is compact on $L^2$ we may assume $Gu_k$ converges; then
$u_k=\tilde I_D\tilde bu_k+Gu_k$ converges in $L^2$ to some $u$, which lies in the
closed subspace $\mathcal{N}^0$, has norm $1$ and satisfies $\tilde bu=0$ by
continuity of $\tilde b$ (on $\mathcal N^0$ the graph norm is the $L^2$ norm),
contradicting injectivity. The estimate gives closedness of the range and continuity
of the inverse.
\end{proof}

Thus $\tilde I_D$ is a left parametrix for $\tilde b$, in the sense that $\tilde I_D
\tilde b=\Id+G$ on $\mathcal{N}^0$ (where $\tilde I_D$ maps into $L^2(X;V_1)\supset
\mathcal{N}^0$).
